\documentclass[11pt,a4paper]{article}
\usepackage[margin=22mm]{geometry}
\usepackage{amsmath,amssymb,amsthm,bm,mathtools}
\usepackage[T1]{fontenc}
\usepackage{lmodern,microtype}
\usepackage{booktabs}
\usepackage{array}
\usepackage{longtable}
\usepackage{enumitem}
\usepackage{float}
\usepackage[htt]{hyphenat}
\emergencystretch=1.5em
\usepackage[colorlinks=true,linkcolor=blue,urlcolor=blue,citecolor=blue]{hyperref}
\newcommand{\R}{\mathbb R}
\newcommand{\E}{\mathbb E}
\newcommand{\tr}{\operatorname{tr}}
\newcommand{\col}{\operatorname{col}}
\newcommand{\sX}{\hat s_X^2}
\newcommand{\sY}{\hat s_Y^2}
\newcommand{\amp}{\mathcal A}
\newcommand{\code}[1]{\texttt{#1}}
\newcommand{\agree}{\textbf{agree}}
\newcommand{\differ}{\textbf{differ}}
\newtheorem{remark}{Remark}
\title{Mapping the linear-solution predictions onto \texttt{fct}:\\discrepancies, conventions and an experiment plan}
\author{}
\date{}
\begin{document}
\maketitle

\begin{abstract}
The note \path{linear_solution_limits.tex} ends with eight numerical predictions for a linear one-hidden-layer network on a ring. This note turns them into a plan that a later session can implement without reopening the theory. It compares, convention by convention, what Sections~1--2 of the notes (after the proofreading pass) define with what the \code{fct} library implements, each claim about the code cited by file and line at one commit (named in Section~\ref{sec:scope}); lists the discrepancies with a concrete code change, the tests that move and the cache impact of each; gives the dictionary from every theory symbol to a \code{Config} field; specifies the measurement (tuning maps, their unitary DFT power, the summary statistics) as functions with inputs, outputs and maths; and writes out every experiment with its \code{fct} invocation, a smoke version that runs on a CPU in seconds, the full version, the predicted numbers and the figure. The four mismatches the limits note already lists (uncentred latents, Gaussian teacher, no activity noise, the rate-penalty coefficient) are confirmed; the forward pass, the loss prefactor, the temperature, the kernel and its spectrum agree exactly (the code's ring kernel equals the script's to the bit). A further discrepancy matters more than any of those for running the experiments: the eigenvalue floor of the penalty matrix $M = K^{-1}$ changes the prior the network is trained under and makes SGLD on a smooth kernel prohibitively stiff, which a whitened parametrisation of the spatial prior removes. The activity noise $\eta$ is treated as a formal regulariser: every prediction run here is stated at $\sigma\to0$, where the three ways of sampling $\eta$ agree, and the two predictions that need $\sigma > 0$ (4 and 5) are deferred to the sampler of \path{input_noise_average.tex}. Each prediction is measured on the final states of the $B$ members of an ensemble, one sample per member. The plan closes with an ordered list of work items in the lab's format, most of which are now implemented on the integration branch \code{sigma0}.
\end{abstract}

\section*{Summary of conclusions}
\addcontentsline{toc}{section}{Summary of conclusions}
\begin{itemize}[nosep]
  \item \emph{Agreements.} Forward normalisations $1/\sqrt m$, $1/\sqrt N$, $1/N$; the data term $\frac N{2Pn}\sum(f-y)^2$ with \code{train.beta} the notes' $\beta$, hence $\kappa = \beta/(Pn)$; the prior $K^{(l)} = K^{\mathrm{post}}\otimes K^{\mathrm{pre}}$ with $M = K^{-1}$ and $\mathrm{Cov}(W) = K/\lambda_w$; the wrapped Gaussian ring kernel and its Fourier eigenvalues (bit-identical to the script's for \code{box} $= N$, \code{length\_scale} $= \ell$); periodic boundaries; the rate penalty on every hidden layer with $r = \code{rate\_lambda}$; the identity activation; one hidden layer; and, for $K = I$, the mean-field solver's saddle-point equations (its tilt matrix is $-B$, no factor of $i$ or $N$).
  \item \emph{Discrepancies, in order of importance for the experiments:} (D1) the penalty parametrisation floors the prior spectrum at \code{min\_rel\_eig}$\,\lambda_0$ and makes SGLD stiff (at $\ell = 4$, $N = 64$ the default floor needs $\mathrm{lr} < 5\times10^{-6}$ at $\beta = 0.25$); (D2) no activity noise $\eta$, which by policy stays absent ($\sigma\to0$; \path{input_noise_average.tex}); (D3) latents on $[-a,a]$ in the pinned code, a Gaussian rather than orthonormal teacher, and a latent-lattice test set built for centred latents; (D4) the mean-field solver knows only the iid prior, so the spatial theory has no implementation in \code{fct}; (D5) naming: \code{net.g} is a hidden-layer gain, unrelated to the notes' rate coefficient $r$; $n = m$ is forced; $k_0$, $k_1$ are kernel amplitudes. On \code{sigma0}, D1 is addressed by \code{spatial.parametrisation = "whitened"} with \code{net.init = "prior"}, D3 by \code{data.latent\_min}, \code{data.latent\_max} and \code{data.teacher = "orthonormal"}, and D4 by \path{fct/linear_theory.py} (the per-mode theory, not the solver).
  \item \emph{Measurement.} A monitor (\code{train.monitor\_kind = "spectrum"}) records, every \code{monitor\_int} SGLD steps, the unitary DFT power of the hidden activity projected on the singular vectors of the latent matrix $Z$, the mean power of the $P - K$ null directions and, at $K = 3$, the hexatic moment $\langle e^{6i\theta}\rangle$ of the tuning angles. One ensemble member gives one sample: its final state. Every predicted statistic (the $q = 0$ share, $R(\Delta)/R(0)$, widths, ratios, $q_{1/2}$) is a statistic of the mean spectrum, so it is read off the final-state spectrum averaged across the $B$ members, with error bars by a bootstrap over members; per-member values serve only as spreads and diagnostics, and the records along a chain only diagnose burn-in. The pure functions are in \path{fct/spectrum.py}; the theory side is \path{fct/linear_theory.py}, a pure function from a \code{Config} to $d^{(k)}_q$ (a port of the per-mode solver of \path{check_linear_soln_limits.py}).
  \item \emph{Experiments.} Predictions 1--3 and 6--8 are runnable from one preset (\code{linear\_ring}) with \code{--set} and \code{--axis} overrides, each once per parametrisation of the spatial prior (\code{"penalty"} and \code{"whitened"}), which must agree on the final-state statistics; a smoke preset (\code{linear\_ring\_tiny}: $N = 16$, $P = 32$, $10^4$ steps, $B = 16$, \code{box} $= N$) runs in seconds and its predicted numbers are given. Every number is at $\sigma^2 = 0$; predictions 4 and 5 are deferred. What each tests: 1 that weak-target maps are prior-shaped, 2 that a rate penalty whitens the strong mode, 3 that temperature condenses it, 4 and 5 the noise floor and its level, 6 that no cell type aligns with an individual latent, 7 that the correlation length is the kernel's, 8 that a common prior scale acts non-monotonically. Order of work: W1--W3 (latent bounds, orthonormal teacher, whitened prior) change conventions; W5--W9 build the measurement (theory from a config, the pure spectrum module, the monitor and its wiring, presets with the end-to-end smoke test, plots); W10--W12 are the runs. On \code{sigma0}, W1--W3 and W5--W7 are done, W8 is in progress and W9 is to do; W4 (activity noise) is withdrawn.
  \item \emph{Caveat on one statistic.} The $q = 0$ share of a condensed mode falls like $1/N$ at fixed $\tau$ (0.343 at $N = 64$, 0.172 at 128, 0.086 at 256): prediction 3's numbers hold at $N = 64$ only; the lag-1 autocorrelation ($0.985$) is $N$-independent.
  \item \emph{Caveat on the sample size.} One sample per member makes the $q = 0$ components of the member-averaged spectrum a mean of $B$ $\chi^2_1$ draws: at $B = 16$ the share of P1 has a standard error of $0.05$. The falsifiers are therefore ring-averaged statistics ($R(\Delta)/R(0)$, the band share of $|q|\le3$, widths), testable at $B\approx16$--$230$; single-$q$ shares, ratios and $q_{1/2}$ are secondary, for which $B$ is raised with \code{exp extend} when needed. $B\ge16$ is the minimum, and the table of Section~\ref{sec:protocol}(ii) sets $B$ per falsifier.
  \item \emph{Caveat on finite $N$.} The per-mode theory is the mean-field limit $N\to\infty$ at fixed $n$, and at finite $N$ it misses an $O(n/N)$ correction (D4): on the full preset ($N = 64$, $n = 8$) at $\beta = 10$ the exact $q = 0$ power is $0.57$ of the predicted. Experiments are therefore run only where $N\gg n$ and the theory is accurate (the regime requirement of W10 and W11).
\end{itemize}

\tableofcontents

\section{What was read, and how to read this note}
\label{sec:scope}
Theory: \path{latex/sections/01_setup.tex} and \path{02_renormalisation_group.tex} after the proofreading pass (their red notes in \path{latex/notes/setup_*.tex}, \path{rg_*.tex}), \path{linear_solution_limits.tex} (the predictions), \path{linear_without_ansatz.tex} and \path{continuum_limit.tex}. The corrected implicit solution is
\begin{equation}
  1 - \frac\delta c = \frac cN\sum_q\frac{h_q}{c^2 - \sY h_q},\quad
  h_q = \frac{\lambda_q(\lambda_q\sX + \sigma^2)}{1 + r(\lambda_q\sX + \sigma^2)},\quad
  d_q = \frac{c^2h_q/\lambda_q}{c^2 - \sY h_q},
  \label{eq:linear-soln}
\end{equation}
on the branch $c > \hat s_Y\sqrt{\max_q h_q}$, with $\sX = k_0 s_X^2/m$, $\sY = s_Y^2/k_1$, $\delta = 1/(k_1\kappa)$, $\kappa = \beta/(Pn)$.

Code: the \code{fct} package and \code{tests/} at commit \code{4f74b73} (the branch \code{latex-edit}). Every \code{file:line} below refers to that commit, and Section~\ref{sec:side} and the \emph{Where} paragraphs of Section~\ref{sec:discrepancies} describe the code as it was there: they are the record of what had to change. At that commit the linear-teacher latents were drawn on $[-a, a]$ and the teacher was Gaussian, $W\sim\mathcal N(0, 1/K)$.

\emph{After the pinned commit.} The integration branch \code{sigma0} (at \code{951153c}) has merged the work items that address the discrepancies; where this note describes code that exists only there it names modules, functions and fields, not lines. In order: the policy $\sigma\to0$ and its note \path{input_noise_average.tex} (D2); \code{latent-bounds} and \code{orthonormal-teacher}: latents uniform on $[\code{data.latent\_min}, \code{data.latent\_max}]$, the lattice test set at $\{\min, \mathrm{mid}, \max\}$, \code{data.teacher = "orthonormal"} (W1, W2; D3); \code{spectrum-module}: \path{fct/spectrum.py} (W6); \code{linear-theory-module}: \path{fct/linear_theory.py} (W5; D4); \code{whitened-spatial-prior}: \code{spatial.parametrisation} and \code{net.init = "prior"} (W3; D1); \code{rate-penalty-latex-r} and \code{rate-penalty-code}: $r = \code{rate\_lambda}$ (\code{LOSS\_VERSION = 5}; Section~\ref{sec:loss}); \code{spectrum-monitor}: \code{train.monitor\_kind = "spectrum"} (W7); \code{ensemble-extend}: \code{fct exp extend ROOT --by K} adds members to every setting of an experiment. The presets (W8, \code{linear-ring-presets}) are in progress and the plots (W9, \code{ring-plots-notebook}) to do.

Notation follows the notes: layer superscripts $(l)$, data indices $\mu,\nu$ superscript, spatial indices $i,j,a,b,\alpha,\beta$ subscript; $K$ is the number of latents where the context is the task and the kernel otherwise, as in the limits note. $r$ is the rate-penalty coefficient, $\Delta$ a distance or lag on the ring ($R(\Delta)$, $K(\Delta)$), $\hat\delta$ the reduced ridge of the limits note; \code{net.g} is a gain unrelated to any of them. $\sigma^2$ appears in the theory's formulas, and is $0$ in every number of this note. Code identifiers are in \code{typewriter}; a dotted path such as \code{spatial.hidden.kernel.length\_scale} is a \code{Config} field.

\section{Architecture and conventions, side by side}
\label{sec:side}

\subsection{Forward pass, normalisations and the input noise}
\label{sec:forward}
\emph{Notes} (\code{eq:ff-network}): $h^{(1)\mu}_i = m^{-1/2}\sum_a W^{(0)}_{i,a}x^\mu_a + \eta^\mu_i$, $h^{(l+1)\mu}_i = N^{-1/2}\sum_j W^{(l)}_{i,j}\phi(h^{(l)\mu}_j)$, $f^\mu_\alpha = N^{-1}\sum_j W^{(L)}_{\alpha,j}\phi(h^{(L)\mu}_j)$, with $\eta^\mu_i\sim\mathcal N(0,\sigma^2)$ i.i.d.\ over $(\mu, i)$, integrated out in \code{eq:single-site-reg} into the term $-\frac{\sigma^2}2[\hat h^{(1)}]^2$, which is why the first-layer covariance is $\Phi^{(0)} + \sigma^2 I$ (\code{02\_renormalisation\_group.tex}, the single-site density at its top).

\emph{Code} (\path{fct/model.py}, class \code{MLP}): \code{forward} at lines 83--96.
\begin{itemize}[nosep]
  \item line 88, \code{(1 / math.sqrt(self.D\_in)) * einsum("id, pd -> pi", self.U, x)}: $h^{(1)} = Ux/\sqrt m$ with $m = $ \code{D\_in};
  \item line 92, \code{(self.g / math.sqrt(self.N)) * einsum("ij, pj -> pi", J, z\_prev)}: the hidden-to-hidden step with a gain \code{g} in front of $1/\sqrt N$ (unused at $L = 1$: \code{J} is empty, line 47);
  \item line 95, \code{(1 / self.N) * einsum("id, pi -> pd", self.V, self.phi(H[-1]))}: the $1/N$ readout of the post-activations (\code{MODEL\_VERSION = 2}, \path{fct/config.py:35}).
\end{itemize} The ensemble class repeats the same three factors (lines 160--165). All parameters are standard-normal at initialisation (lines 42--47); the scale factors live in the forward pass, which is the notes' convention of unit-variance weights. \emph{The forward pass is deterministic: there is no $\eta$} (nothing is added to \code{H[0]}; no RNG is read in \code{forward}).

\emph{Verdict:} $1/\sqrt m$, $1/\sqrt N$ (with \code{g = 1}), $1/N$: \agree. Input noise $\eta$: \differ{}, and kept so by policy ($\sigma\to0$; D2, Section~\ref{sec:D2}). Tests pinning the normalisations: \path{tests/test_model.py:46} and \code{:53}.

\subsection{Loss prefactor, temperature, SGLD, and \texorpdfstring{$\kappa$}{kappa}}
\label{sec:loss}
\emph{Notes:} $\mathcal L = \frac N{2Pn}\sum_{\mu,\alpha}(f^\mu_\alpha - y^\mu_\alpha)^2 + \frac1\beta\mathcal P_{\mathrm{weight}} + \frac1\beta\mathcal P_{\mathrm{rates}}$, Langevin $\partial_s\Theta = -\nabla\mathcal L + \sqrt{2/\beta}\,\eta(s)$, stationary law $e^{-\beta\mathcal L}$ (\code{eq:theta-distr}); $\kappa = \beta/(Pn)$ is defined where $f$ is integrated out.

\emph{Code:} \path{fct/train.py:57--95}, \code{compute\_loss}: line 85 \code{mse = per\_member(net, (y - z).square(), torch.mean)}, the mean over the $P\times n$ entries; line 93 \code{loss = 0.5 * net.N * mse + prior\_coefficient * weight\_l2 + rate\_lambda * rate\_l2}. So $\frac N2\,\mathrm{mse} = \frac N{2Pn}\sum(f-y)^2$ exactly, with $n = $ \code{D\_out}. \path{fct/sgld.py:53--76}, \code{SGLD.step}: line 60 \code{noise\_scale = math.sqrt(2 * group["lr"] / group["beta"])}, line 64 \code{delta = noise\_scale * self.noise(p) - group["lr"] * p.grad}: the Euler--Maruyama step of the notes' Langevin equation with $\mathrm dt = $ \code{lr} and $\beta = $ \code{train.beta} (\path{fct/config.py:121}). The stationary law is $e^{-\beta L}$ as $\mathrm{lr}\to0$; at finite \code{lr} the variance on a quadratic of curvature $c$ is $1/(\beta c(1 - \mathrm{lr}\,c/2))$ (\path{docs/design.md}, section ``SGLD''; \path{tests/test_sgld.py:45}). The prior coefficient (\path{fct/train.py:45--54}, \code{prior\_coefficient}) is $\lambda_w/(2\beta)$ so that $\beta L$ carries $\frac{\lambda_w}2\,\mathrm{form}$, i.e.\ the notes' $\mathcal P_{\mathrm{weight}} = \frac12 W M W$ with $M = \lambda_w(M^{\mathrm{post}}\otimes M^{\mathrm{pre}})$ (\code{LOSS\_VERSION = 4}, \path{fct/config.py:29}). At the commit above, line 93 adds the rate term unscaled, $\code{rate\_lambda}\sum\phi^2$. \code{LOSS\_VERSION = 5}, which postdates that commit, scales the rate term by $1/(2\beta)$ (\code{rate\_coefficient}, the rule of the prior coefficient), so $\beta L$ carries $\frac{\code{rate\_lambda}}2\sum\phi^2$ against the notes' $\frac r2\sum\phi^2$ and $r = \code{rate\_lambda}$, the solver's \code{rate\_tilt} (\code{network\_constants}).

\emph{Verdict:} \agree{} with \code{train.beta} $=\beta$, \code{train.optimizer = "sgld"}, $\kappa = \beta/(Pn)$ and $\code{rate\_lambda} = r$ (from \code{LOSS\_VERSION = 5}; at the pinned commit $r = 2\beta\,\code{rate\_lambda}$). The mean-field solver reads the same $\beta$ through \code{resolve\_beta} (\path{fct/meanfield.py:277--282}). Under Adam there is no Gibbs law; every experiment here uses SGLD.

\subsection{The weight prior}
\label{sec:prior}
\emph{Notes:} $\mathcal P_{\mathrm{weight}} = \frac12\sum_l W^{(l)}_{i,a}M^{(l)}_{i,a,j,b}W^{(l)}_{j,b}$, $K = M^{-1}$, $\E[W_{i,a}W_{j,b}] = K_{i,a,j,b}$, factorised as $K^{(l)}_{i,a,j,b} = K^{(l),\mathrm{post}}_{i,j}K^{(l),\mathrm{pre}}_{a,b}$; Section~2 sets $K^{(0),\mathrm{post}} = K^{(1),\mathrm{pre}}\equiv K$, $K^{(0),\mathrm{pre}} = k_0I_m$, $K^{(1),\mathrm{post}} = k_1I_n$, $K = F\Lambda F^\dagger$ with the unitary DFT (translation invariance, periodic boundaries), $\lambda_q$ real, $\lambda_q = \lambda_{-q}$; the limits note fixes $K_{i,i} = 1$ (so $\bar\lambda = 1$) and the wrapped Gaussian of \path{fct/spatial.py}.

\emph{Code:} \path{fct/spatial.py}. The prior of layer $l$ is $K^{\mathrm{post}}\otimes K^{\mathrm{pre}}$ per population pair (module docstring lines 1--15; \code{LayerSpace} 319--334; \code{build\_spaces} 354--363, which gives layer $l$ the populations $l$ and $l + 1$ and shares one \code{Population} object between consecutive layers, line 363, so $K^{(0),\mathrm{post}} = K^{(1),\mathrm{pre}}$ holds by construction). The penalty is per factor, $M = K^{-1}$ by \code{penalty\_matrix} (269--297), and the loss uses $\tr(M^{\mathrm{post}}WM^{\mathrm{pre}}W^\top)$ without forming the Kronecker product (\code{weight\_quadratic\_form}, 369--380; \code{parameter\_factors}, 435--447, with the readout \code{V} stored $(N, n)$ so its rows take $M^{\mathrm{pre}}$). Combined with the prior coefficient of Section~\ref{sec:loss}, $\mathrm{Cov}(\mathrm{vec}\,W^{(l)}) = (K^{\mathrm{post}}\otimes K^{\mathrm{pre}})/\lambda_w$ (\path{docs/design.md}, ``Spatial weight prior in training''; \path{tests/test_prior_convention.py:46}). The kernel: \code{periodic\_gaussian\_kernel} (224--231) is the product over axes of \code{wrapped\_gaussian\_profile} (208--221), $K_{i,j} = \code{amplitude}\cdot W(\delta)/W(0)$ with $W(\delta) = \sum_{|k|\le n_{\mathrm{img}}}\exp(-(\delta + k\,\code{box})^2/2\ell^2)$, $n_{\mathrm{img}} = \lceil 8\ell/\code{box}\rceil + 1$ (201--205); the diagonal is \code{amplitude} exactly. A grid embedding puts $n$ units at $\code{box}\cdot k/n$, $k = 0,\dots,n-1$ (\code{\_grid\_positions}, 160--163), on the torus $[0,\code{box})$ (\code{torus\_distance}, 195--198; the reduction at 217). \code{length\_scale = 0} gives \code{amplitude}$\cdot I$ (218--219). Scalars: $k_0$ and $k_1$ are \code{spatial.input.kernel.amplitude} and \code{spatial.output.kernel.amplitude} with \code{length\_scale = 0} on both and \code{weight\_lambda = 1}; the hidden amplitude is $1$ ($K_{i,i} = 1$). \emph{Checked numerically} (Appendix~\ref{app:numbers}): with $N = 64$, \code{box = 64.0}, grid embedding, hidden \code{length\_scale = 4.0}, \code{amplitude = 1.0}, the hidden-population \code{K} of \code{build\_spaces} equals the script's \code{ring\_kernel(64, 4.0)} with \code{max|diff| = 0.0}, as does \code{box = 1.0} with \code{length\_scale = 4/64}; its eigenvalues equal \code{ring\_spectrum(64, 4.0)} to $9\times10^{-15}$, $\lambda_0 = 10.0265$, $\bar\lambda = 1$. So the limits note's $\ell$ in sites is \code{length\_scale} $=\ell$ with \code{box} $= N$ (or $\ell/N$ with \code{box} $= 1$), and the Fourier eigenvalues, the periodicity and $\bar\lambda = 1$ all \agree.

The input/output positions $\mathbf x_a$, $\mathbf x_\alpha$: the notes assign them but with $K^{(0),\mathrm{pre}} = k_0I$, $K^{(1),\mathrm{post}} = k_1I$ nothing depends on them; the code assigns grid positions to the input and output populations too (\code{build\_spaces}, 360 and 362) and with \code{length\_scale = 0} they drop out identically. \agree.

What does \emph{not} agree is the spectrum the network is actually trained under. \code{penalty\_matrix} floors the eigenvalues of $K$ at \code{min\_rel\_eig}$\cdot\lambda_0$ before inverting (lines 285--287), so the sampled prior is $K_{\mathrm{eff}} = M^{-1}$ with $\lambda_q^{\mathrm{eff}} = \max(\lambda_q, f\lambda_0)$, $f = $ \code{spatial.min\_rel\_eig} (default $10^{-6}$, line 118; the presets use $10^{-2}$). The docstring (41--46) and \path{docs/design.md} (``The floor'') say so. This is D1 (Section~\ref{sec:D1}): \differ.

\subsection{The firing-rate regulariser}
\emph{Notes:} $\mathcal P_{\mathrm{rates}} = \frac12\sum_{l=1}^L\sum_\mu\sum_{i,j}\phi^{(l)\mu}_iG^{(l)}_{i,j}\phi^{(l)\mu}_j$, $G = rI$: every hidden layer, not the output.

\emph{Code:} \path{fct/train.py:86}, \code{rate\_l2 = per\_member(net, net.phi(h).square())} with \code{h} the \code{(P, L, N)} stack of hidden pre-activations (\code{MLP.forward} returns \code{torch.stack(H, dim=1)}, line 96, hidden layers only); summed over patterns, neurons and layers (\path{tests/test_train.py:269}). Coefficient: $r = \code{rate\_lambda}$ (Section~\ref{sec:loss}). \agree{} in layer range and form. The code's \code{net.g} (\path{fct/config.py:93}) is the gain of the hidden-to-hidden weights (Section~\ref{sec:forward}), unrelated to $r$, and must stay $1$ (D5).

\subsection{Activation and depth}
\code{net.phi = "id"} maps to \code{fct.activations.identity} (\path{fct/activations.py:11,21}) with derivative one and tilt class \code{"linear"} (lines 65, 79); \code{net.L = 1} builds \code{U} $(N, m)$ and \code{V} $(N, n)$ and no \code{J} (\path{fct/model.py:46--47}). \agree{} (\path{tests/test_identity_activation.py:38,51}).

\subsection{The task}
\label{sec:task}
\emph{Notes} (limits note, Section~2.5): $X = ZW_X^\top$, $Y = ZW_Y^\top$, $Z\in\R^{P\times K}$ with $z^\mu_k$ i.i.d.\ uniform on $[0,a]$, $W_X\in\R^{m\times K}$, $W_Y\in\R^{n\times K}$ with orthonormal columns; then $XX^\top = YY^\top = ZZ^\top$, the ansatz is exact, $s_1^2 = Pa^2(3K+1)/12$ on the mean direction $u_1\propto Z\mathbf 1$ and $s_{k\ge2}^2 = Pa^2/12$ on the contrasts.

\emph{Code} (\path{fct/tasks.py:81--103}, \code{linear\_teacher}): line 90 \code{z = a * (2 * torch.rand((cfg.P, k), ...) - 1)}: $z\sim\mathcal U[-a, a]$; lines 91--92 \code{w\_x}, \code{w\_y} \code{= torch.randn((m, k), ...) / k**0.5}: i.i.d.\ $\mathcal N(0, 1/K)$ entries, not orthonormal columns; \code{m = cfg.D\_in} for both $x$ and $y$ (line 88; \code{output\_dim = lambda cfg: cfg.D\_in}, line 264), so $n = m$ is forced; $K = $ \code{data.n\_latents} and $a$ the half-width field of \code{DataConfig} (\path{fct/config.py:60--61}); \code{factors} hold the columns of $Z$ (line 100) and \code{extras} the maps. The held-out lattice \code{latent\_lattice\_dataset} (353--385) uses \code{values = [-a, 0, a]} (line 374), the three-point lattice of a \emph{centred} range. Tests that pin the current conventions: \path{tests/test_tasks.py:158} (zero mean, variance $4/3$ at $a = 2$), \code{:185} (teacher variance $1/K$), \path{tests/test_latent_lattice.py:26,48}.

\emph{Verdict:} \differ{} on the range, on the teacher, and on the lattice; $m = n$ is a restriction the predictions satisfy ($m = n = 8$). D3 (Section~\ref{sec:D3}), resolved on \code{sigma0}: \code{data.latent\_min}, \code{data.latent\_max} (uniform on $[\min, \max]$; the presets use $[0, 1]$, the notes' $[0, a]$ with $a = 1$) and \code{data.teacher = "orthonormal"}.

\subsection{What the mean-field solver solves}
\label{sec:solver}
\emph{Notes:} the action of Section~1 after the contour rotation $\hat\Phi\to-i\hat\Phi$ (red note, \path{setup_conjugate_rotation.tex}): single-site density $\propto\exp\{-\frac12h^\mu_i(\Phi^{-1})^{\mu,\nu}_{i,j}h^\nu_j + \frac12\phi^\mu_i(\hat\Phi^{\mu,\nu}_{\alpha,\beta}K_{\alpha,i,\beta,j} - G_{i,j}\delta^{\mu,\nu})\phi^\nu_j\}$, kernel $\Phi^{(l)} = \frac1N\sum_{a,b}K^{(l)}_{i,a,j,b}\langle\phi_a\phi_b\rangle$, readout saddle $\hat\Phi^{(L)} = (\Phi^{(L)} + \kappa^{-1}I)^{-1}\mathbb Y(\Phi^{(L)} + \kappa^{-1}I)^{-1}$, hidden-layer saddle \code{eq:backward-pass} $\hat\Phi^{(l)} = \frac1N([\Phi^{(l)}]^{-1}\langle h\otimes h\rangle[\Phi^{(l)}]^{-1} - [\Phi^{(l)}]^{-1})$. In Section~2 with $K = I$ and $k_0 = k_1 = k$: per-site prior $\Phi^{(0)} = \frac km XX^\top + \sigma^2I$, tilt $+\frac12\phi B\phi - \frac{r}2\phi\phi$ with $B = \frac1k\tilde C^{-1}YY^\top\tilde C^{-1}$, $\tilde C = C + \delta I$, $\delta = Pn/(\beta k)$.

\emph{Code:} \path{fct/meanfield.py}, module docstring lines 5--28 and \path{fct/hmc.py:7--16}: the single-site density of layer $l$ is $p(h)\propto\mathcal N(h; 0, \Sigma)\exp(-\frac12\psi^\top M\psi - \frac12\lambda_r\phi^\top\phi)$ with $\Sigma_0 = \code{prior\_scale}\cdot XX^\top/D_{\mathrm{in}}$ (\code{\_prior\_cov}, 285--287; \code{Phi0 = X @ X.T / D\_in}, line 253), $M = \hat\Phi_{\mathrm{code}}$, and the readout target $M_L = -2c\,A\,YY^\top A$, $A = (\alpha I + \Phi_L)^{-1}$ (\code{\_last\_stein\_target}, 682--684; \code{\_learning\_pieces}, 687--695). The constants follow the network (\code{network\_constants}, 263--274): \code{prior\_scale} $= 1/\lambda_w$, $\alpha = \lambda_wPn/\beta$, \code{rate\_tilt} $= \code{rate\_lambda}$, $c = \lambda_w/2$. With $k = 1/\lambda_w$: $\Sigma_0 = \frac km XX^\top$ ($\sigma^2 = 0$), $\alpha = Pn/(\beta k) = \delta$, $-\frac12\psi M_L\psi = +\frac12\psi\,\lambda_w(\tilde C^{-1}YY^\top\tilde C^{-1})\psi = +\frac12\psi B\psi$, and $\lambda_r = r$. The hidden-layer gradient (\code{\_outer\_gradients}, 945--974, lines 961--973) is $-\frac12\hat\Phi_l + \frac12\,\mathrm{sym}(\langle uh^\top\rangle_{l+1}\Phi_l^{-1})$, the Stein form of \code{eq:backward-pass} with the opposite sign. The kernels are $\Phi_l = \psi\psi^\top/N$ (\code{empirical\_kernels}, 153--161), the notes' $C$ for $K = I$. The analytic reference of \path{tests/test_meanfield.py:77} iterates exactly $\Phi = (\Sigma^{-1} + \hat\Phi)^{-1}$, $\hat\Phi = -2cAYY^\top A$.

\emph{Verdict:} for $K = I$ the solver's saddle-point equations are the notes' with the dictionary $M_{\mathrm{code}} = -B = -K^{(1),\mathrm{post}}_{\alpha,\beta}\hat\Phi^{(1)\mu,\nu}_{\alpha,\beta}$ (the code's tilt is a penalty, the notes' a reward; no factor of $i$ on either side, both being the rotated conjugate; no factor of $N$, the per-site tilt is $\frac12\phi B\phi$ in both). \agree{} on the iid prior. For a spatial prior the solver refuses to run (\code{SpatialPriorUnsupportedError}, \path{fct/meanfield.py:1513--1516}): the spatial theory of Section~2 has no implementation. D4 (Section~\ref{sec:D4}).

\subsection{Summary table}
\begin{table}[H]
\centering\footnotesize
\begin{tabular}{@{}>{\raggedright\arraybackslash}p{3.1cm}>{\raggedright\arraybackslash}p{3.6cm}>{\raggedright\arraybackslash}p{5.4cm}>{\raggedright\arraybackslash}p{2.6cm}@{}}
\toprule
Item & Notes & \code{fct} (the commit of Section~\ref{sec:scope}) & \\
\midrule
$h^{(1)} = Ux/\sqrt m$ & \code{eq:ff-network} & \path{model.py:88} & agree\\
$f = V^\top\phi(h)/N$ & \code{eq:ff-network} & \path{model.py:95} & agree\\
$\eta\sim\mathcal N(0,\sigma^2)$ & \code{eq:ff-network}, \code{eq:single-site-reg} & absent; $\sigma\to0$ by policy & D2\\
data term $\frac N{2Pn}\sum(f-y)^2$ & $\mathcal L$ & \path{train.py:85,93} & agree\\
$\beta$, SGLD, $\kappa = \beta/(Pn)$ & \code{eq:theta-distr} & \path{sgld.py:60,64}; \code{train.beta} & agree\\
$\mathcal P_{\mathrm{weight}} = \frac12WMW$, $K = M^{-1}$ & \code{eq:prior-weights} & \path{train.py:45--54,104--110}; $\mathrm{Cov} = K/\lambda_w$ & agree\\
$K^{(l)} = K^{\mathrm{post}}\otimes K^{\mathrm{pre}}$ & Section~1 & \path{spatial.py:319--363} & agree\\
$k_0$, $k_1$ & Section~2 & input/output \code{amplitude}, $\ell = 0$, $\lambda_w = 1$ & agree (dictionary)\\
$\lambda_q$, periodic, $K_{i,i} = 1$ & \code{eq:ring} & \path{spatial.py:201--231}; bit-identical & agree\\
spectrum sampled & $\lambda_q$ & $\max(\lambda_q, f\lambda_0)$, \path{spatial.py:285--287} & D1\\
$G = rI$, layers $1..L$ & $\mathcal P_{\mathrm{rates}}$ & \path{train.py:86}; $r = $ \code{rate\_lambda} & agree\\
$\phi = \mathrm{id}$, $L = 1$ & Section~2 & \code{net.phi="id"}, \code{net.L=1} & agree\\
$Z\sim\mathcal U[0,a]$ & limits note & $\mathcal U[-a,a]$ (\path{tasks.py:90}); \code{sigma0}: $[\min,\max]$ & D3\\
$W_X$, $W_Y$ orthonormal & limits note & $\mathcal N(0,1/K)$ (\path{tasks.py:91--92}); \code{sigma0}: \code{teacher} & D3\\
lattice test set & -- & $\{-a,0,a\}^K$ (\path{tasks.py:374}); \code{sigma0}: $\{\min,\mathrm{mid},\max\}^K$ & D3\\
saddle point, $K = I$ & \code{eq:backward-pass} & \path{meanfield.py:682--695,961--973}, $M = -B$ & agree\\
saddle point, spatial & Section~2 & refused (\path{meanfield.py:1513}) & D4\\
\bottomrule
\end{tabular}
\caption{Conventions side by side. ``agree'' means equal under the dictionary of Section~\ref{sec:dictionary}.}
\end{table}

\section{Discrepancies and how to address them in the code}
\label{sec:discrepancies}
Ordered by how much they matter for the experiments. Each entry: where, what the notes say, the change, the tests that move, the cache impact (\code{Config.section\_hash}, \path{fct/config.py:391--426}).

\subsection{D1: the eigenvalue floor of the penalty, and SGLD stiffness}
\label{sec:D1}
\emph{Where.} \path{fct/spatial.py:269--297} (\code{penalty\_matrix}), used by \code{\_make\_population} (337--351) for every population; the stability check \path{fct/train.py:113--136}.

\emph{Notes.} The prior spectrum is $\lambda_q$ of \code{eq:ring} (\path{linear_solution_limits.tex:264}), which for $\ell = 4$, $N = 64$ spans $\lambda_q/\lambda_0 = 1$ down to $10^{-34}$ at $|q| = 32$.

\emph{What the code does.} $M = K^{-1}$ needs the inverse, so eigenvalues below $f\lambda_0$ are raised to the floor and the sampled prior has $\lambda_q^{\mathrm{eff}} = \max(\lambda_q, f\lambda_0)$. Measured with \code{build\_spaces} at $N = 64$, $\ell = 4$ (Appendix~\ref{app:numbers}): $f = 10^{-2}$ floors 49 of 64 eigenvalues (every $|q|\ge8$, floor $0.1$), $f = 10^{-3}$ floors 45, $f = 10^{-4}$ floors 43 ($|q|\ge11$, floor $10^{-3}$), $f = 10^{-6}$ floors 37 ($|q|\ge14$). The effect on the $\sigma^2 = 0$ predictions (computed with \path{fct/linear_theory.py} on $\lambda_q^{\mathrm{eff}}$; Appendix~\ref{app:numbers}): the ratio $\hat d^{(1)}_q/\hat d^{(2)}_q$ of P1 is $10$ at every floored $|q|$, as without the floor (at $\sigma = 0$ both modes sit on the same floored $\lambda_q$, scaled by $\hat s_{X,k}^2$), so P1's ratio cannot see the floor, but every shape statistic can. At $f = 10^{-2}$ P1's $R(1)/R(0)$ falls from $0.969$ to $0.881$ and its share from $0.157$ to $0.146$, P2's mode-1 $R(1)/R(0)$ from $0.922$ to $0.371$ (its ratio at $|q| = 12$ is $8.10$ instead of $10.0$), P3's from $0.985$ to $0.951$ (mode 1) and $0.971$ to $0.888$ (mode 2). At $f = 10^{-3}$ P2's mode-1 $R(1)/R(0)$ is $0.812$ and its share $0.071$ (against $0.077$), P1's $R(1)/R(0)$ $0.960$, P3's within $0.009$. At $f = 10^{-4}$ P2's mode-1 $R(1)/R(0)$ is $0.910$ (off by $0.012$) and P1 and P3 are within $0.001$; at $f = 10^{-6}$ nothing moves by more than $10^{-3}$. The price of a small floor is SGLD stiffness: \code{\_check\_sgld\_step\_resolves\_the\_prior} requires $\mathrm{lr}\cdot\frac{\lambda_w}{2\beta}\lambda_{\max}(M^{\mathrm{post}})\lambda_{\max}(M^{\mathrm{pre}}) < 1$ with $\lambda_{\max}(M^{\mathrm{hidden}}) = 1/(f\lambda_0)$ (measured $9.97$, $997$, $9.97\times10^4$ for $f = 10^{-2}, 10^{-4}, 10^{-6}$), so at $\beta = 0.25$ the bound is $\mathrm{lr} < 0.050$, $5\times10^{-3}$, $5\times10^{-4}$, $5\times10^{-6}$ for $f = 10^{-2}, 10^{-3}, 10^{-4}, 10^{-6}$. That bound is the Euler--Maruyama divergence edge of the stiffest prior mode ($\mathrm{lr}\,c < 2$, $c = \lambda_w/(\beta f\lambda_0)$), not a resolution criterion; Section~\ref{sec:protocol}(v). The softest prior mode (variance $\lambda_0k_0$) relaxes in $\beta\lambda_0k_0/\mathrm{lr}$ steps, i.e.\ at the bound $0.5/f$ steps: $50$, $500$, $5\times10^3$, $5\times10^5$. The choice of $f$ trades the floor's bias (above all on P2's mode 1) against the penalty path's step and run length. The real runs use $f = 10^{-6}$, where the floor moves no prediction by more than $10^{-3}$, at the cost of a penalty step below $5\times10^{-6}$ at $\beta = 0.25$ and correspondingly long penalty runs. The tiny smoke preset may use $f = 10^{-3}$ as a pipeline check against its own floored prediction (\code{theory\_parameters} uses the floored $\lambda_q$ of a penalty run).

\emph{Change: a whitened parametrisation of the spatial prior.} Add \code{spatial.parametrisation: Literal["penalty", "whitened"] = "penalty"} (\path{fct/spatial.py}, \code{SpatialConfig}). Under \code{"whitened"} the trained parameters are $\xi^{(l)}$ with $W^{(l)} = S^{\mathrm{post}}\xi^{(l)}S^{\mathrm{pre}}$, $S = K^{1/2}$ the symmetric PSD square root from \code{torch.linalg.eigh} with eigenvalues clamped at $0$ (no inversion, no floor), and the penalty is $\|\xi^{(l)}\|_F^2$ with the unchanged coefficient $\lambda_w/(2\beta)$, so $\mathrm{Cov}(\mathrm{vec}\,W) = (K^{\mathrm{post}}\otimes K^{\mathrm{pre}})/\lambda_w$ \emph{exactly}, for every spectrum, and the curvature of the prior term is $\lambda_w/\beta$ in every direction: the stability bound becomes $\mathrm{lr}\,\lambda_w/\beta < 2$ and every prior mode relaxes in $\beta/(\lambda_w\,\mathrm{lr})$ steps. Interfaces: \code{Population} gains \code{S} (the square root; \code{M} may stay for readers); \code{parameter\_factors} gains a sibling \code{whitening\_factors(spaces) -> \{name: (S\_rows, S\_cols)\}} in the stored orientation of each parameter (for \code{V}, stored $(N, n)$: rows $S^{\mathrm{pre}}$, columns $S^{\mathrm{post}}$, as for \code{M}); \code{MLP} and \code{EnsembleMLP} accept an optional \code{whitening} mapping of registered buffers and apply $W = S_{\mathrm{rows}}\,\xi\,S_{\mathrm{cols}}$ in \code{forward} before the existing scale factors; a pure \code{effective\_weights(net) -> \{name: W\}} returns the weights a reader should analyse (\code{net.U} is $\xi$ under whitening); \code{make\_weight\_penalty} returns the plain \code{weight\_l2} (of $\xi$) under whitening. The SGLD stability check is \emph{kept and extended}, not skipped: the prior-only bound of \code{\_check\_sgld\_step\_resolves\_the\_prior} misses the data curvature (Section~\ref{sec:protocol}(i)), so under whitening \code{make\_weight\_penalty} (or \code{train}, before its first step) estimates $\lambda_{\max}(\nabla^2L)$ at the initial net on the dataset by Lanczos iteration on Hessian-vector products through \code{torch.autograd.grad} (\code{largest\_hessian\_eigenvalue}, at most \code{HESSIAN\_MAX\_ITERATIONS} $= 300$ iterations, relative tolerance \code{HESSIAN\_RTOL} $= 10^{-8}$) and raises, naming \code{lr} and $\lambda_{\max}$, if $\mathrm{lr}\,\lambda_{\max}\ge1$ (half the divergence limit; the rule \eqref{eq:lr-rule} aims at $0.1$), and logs the value otherwise. Under whitening the initial $\xi$ are drawn from their prior, $\mathcal N(0, I/\lambda_w)$ (\code{make\_net} scales its standard-normal draw by $\lambda_w^{-1/2}$; \path{docs/design.md} ``Not covered'' records that the current initialisation is not from the prior), so directions the data does not touch start in equilibrium. Callers that pass \code{net.U}, \code{net.V} as weights must pass \code{effective\_weights} instead: the notebook callers of \code{order\_weights} (\path{fct/spatial_view.py:176}) and \code{ordered\_kernel\_comparison} (\code{:268}), and \code{weight\_kernels} (\path{fct/phase_view.py:364}). Cost: two extra $(N,N)\times(N,m)$ products per step; on the four-thread CPU run of Appendix~\ref{app:numbers} $0.4$\,ms per loss-and-gradient against $2.0$\,ms for the penalty path (a reviewer measured $0.22$ against $0.43$\,ms on other hardware: the ratio is hardware-dependent, the direction is not).

\emph{Tests.} Existing tests are untouched under the default (\path{tests/test_spatial_prior.py}, \path{tests/test_prior_convention.py}, \path{tests/test_train_pins.py} pin the penalty path bit for bit). New: (i) \code{whitening\_factors} with identity kernels reproduces the plain net bit for bit; (ii) for a whitened net with no data ($P = 0$ or $\beta\to0$ is awkward; instead sample $\xi\sim\mathcal N(0, I/\lambda_w)$ directly) the sample covariance of $\mathrm{vec}\,W$ over $10^4$ draws is $(K^{\mathrm{post}}\otimes K^{\mathrm{pre}})/\lambda_w$ to $5\%$ on a $N = 8$ ring; (iii) \code{effective\_weights} equals $S\xi S$; (iv) the smoke experiment of Section~\ref{sec:smoke} as an end-to-end test.

\emph{Cache.} The field enters the train hash only when set to \code{"whitened"} (the pattern of \code{lambda\_max}, \path{fct/config.py:419--424}), so every existing hash is unchanged; a whitened run has a new hash. Presets unchanged.

\emph{As implemented on \code{sigma0}} (\code{whitened-spatial-prior}). The parametrisation is as above, with one change to the initial law: \code{net.init: Literal["standard", "prior"]} (hashed only as \code{"prior"}; a shared field of \code{ExpConfig}, \code{--set init=prior}) draws the trained parameters from the prior they are trained under, $\xi\sim\mathcal N(0, I/\lambda_w)$ under \code{"whitened"} and $W\sim\mathcal N(0, (K^{\mathrm{post}}_{\mathrm{eff}}\otimes K^{\mathrm{pre}}_{\mathrm{eff}})/\lambda_w)$, the floored prior, under \code{"penalty"}. \code{"whitened"} requires \code{"prior"}. Both paths therefore start in equilibrium in every direction the data does not touch, so neither carries a head start. The stationary laws of the two paths agree up to the floor (which at $f = 10^{-6}$ moves no predicted number by more than $10^{-3}$, above), but the dynamics do not: SGLD in $\xi$ is Langevin preconditioned by $S^2$ in $W$, so one \code{train.lr} gives the two paths different numbers of relaxation times, and the penalty path's run-start check forces $\mathrm{lr} < 2\beta f\lambda_0/\lambda_w$ ($5\times10^{-6}$ at $\beta = 0.25$, $f = 10^{-6}$). Running both and comparing them is the protocol of Section~\ref{sec:protocol}(v).

\subsection{D2: no activity noise, and the policy \texorpdfstring{$\sigma\to0$}{sigma -> 0}}
\label{sec:D2}
\emph{Where.} \path{fct/model.py:88}: nothing is added to the first layer.

\emph{Notes.} $\eta^\mu_i\sim\mathcal N(0,\sigma^2)$ enters \code{eq:ff-network} and is integrated together with the weights (\code{eq:single-site-reg}).

\emph{Decision.} No change to the code. The note \path{input_noise_average.tex} shows that the notes' measure over $(\Theta, \eta)$ (the annealed average) differs from both a fixed $\eta$ (quenched) and an $\eta$ redrawn at every SGLD step (the averaged-loss law) at $O(\sigma^2)$, most visibly on the null directions, and that two samplers reproduce it exactly. Its policy is the one this plan follows: $\sigma$ is a formal regulariser, every prediction is stated at $\sigma\to0^+$ where the three measures coincide, and \code{fct} has no $\eta$. So: no \code{Config} field for $\sigma^2$, no noise in the forward pass or the monitor, the predicted numbers of predictions 1--3, 7, 8 are their $\sigma^2 = 0$ values, and predictions 4 and 5, which measure the noise floor $\sigma_0^2 = \sigma^2/(1 + r\sigma^2)$, are deferred until one of that note's samplers is implemented (Sections~\ref{sec:P4}, \ref{sec:P5}). At $\sigma = 0$ the activity of a linear network lies in $\col(Z)$, so the null row of the spectrum is zero up to round-off: a check of the projection, not of the theory. The theory functions keep $\sigma^2$ as an argument with default $0$ (\code{sig2} of \path{fct/linear_theory.py}).

\subsection{D3: the task generator and the lattice test set}
\label{sec:D3}
\emph{Where.} \path{fct/tasks.py:90} (range), \code{:91--92} (teacher), \code{:374} (lattice); \path{fct/data.py:177} hashes the data section only (\code{section\_hash("data")}): there is no analogue of \code{LOSS\_VERSION}/\code{MODEL\_VERSION} for the generator, so a change of the generator under unchanged field names would leave every cached \code{data/dataset.pt}, \code{train/final.pt} and \code{meanfield/final.pt} with a matching hash, silently reused.

\emph{Notes.} $z\sim\mathcal U[0,a]$; orthonormal $W_X$, $W_Y$ (needed for the ansatz to be exact; with the Gaussian teacher $G_X = W_X^\top W_X\ne I$ and the restricted Gram matrices commute only if $G_X(Z^\top Z)G_Y$ is symmetric, \path{linear_without_ansatz.tex}, eq.~restricted-commute; the error was $0.2\%$ on $c$ for a mild skew). The mean direction is the strong mode ($s_1^2/s_k^2 = 3K+1$), which is what predictions 1--3 and 8 rely on.

\emph{Changes, as implemented on \code{sigma0}} (W1 \code{latent-bounds}, W2 \code{orthonormal-teacher}). (a) Both ends of the latent interval are fields, \code{data.latent\_min} and \code{data.latent\_max} (defaults $-1$, $1$, which reproduce the old generator bit for bit; validated \code{latent\_min < latent\_max}); the latents are uniform on $[\code{latent\_min}, \code{latent\_max}]$ and the presets of this plan use $[0, 1]$. The old field is rejected with a message naming the two new ones. No data version is needed: the field names change the hashed dump of the data section, so no cache made under $[-a, a]$ matches a new config. (b) \code{data.teacher: Literal["gaussian", "orthonormal"] = "gaussian"}; \code{"orthonormal"} takes the \code{Q} factor of the same Gaussian draws (\code{torch.linalg.qr(..., mode="reduced")}), so the random stream and every other field are unchanged, and the field is dropped from the hashed dump at \code{"gaussian"}. (c) The lattice test set is $\{\code{latent\_min}, \frac12(\code{latent\_min} + \code{latent\_max}), \code{latent\_max}\}^K$. (d) \code{D\_out = D\_in} stays (the predictions use $m = n$). (e) Symmetric latents have no mean direction (all $K$ singular values of $Z$ degenerate up to $O(P^{-1/2})$), so \code{fct.spectrum.latent\_basis} refuses them, and an experiment with \code{monitor\_kind = "spectrum"} is validated at load: every setting must be \code{linear\_teacher} with latents \code{latent\_basis} accepts (Section~\ref{sec:measure}).

\subsection{D4: no implementation of the spatial saddle point}
\label{sec:D4}
\emph{Where.} \path{fct/meanfield.py:1513--1516}; the two pure solvers exist only as scripts, \path{latex/notes/scripts/check_linear_soln_limits.py} (per-mode scalar equation, functions \code{ring\_spectrum}, \code{h\_of}, \code{solve\_c}, \code{spectrum\_d}, \code{latent\_modes}, \code{half\_power\_q}, \code{autocorr\_lag1}, lines 81--193 and 661--701) and \path{latex/notes/scripts/linear_without_ansatz.py} (the $P\times P$ ansatz-free map, eq.~master of that note).

\emph{Change, as implemented on \code{sigma0}} (W5, \code{linear-theory-module}). A pure module \path{fct/linear_theory.py}: \code{theory\_parameters(cfg, dataset, sig2=0.0, exact\_kernel=False) -> TheoryParameters} (the dictionary of Section~\ref{sec:dictionary}: $\lambda_q$ of the prior the run trains with, i.e.\ the floored spectrum under \code{"penalty"} and $K$'s own under \code{"whitened"} or with \code{exact\_kernel=True}; $\sX$, $\sY$ per mode from the SVD of $Z$ built from \code{dataset.factors}; $\delta$, $r$ from the config; $\sigma^2$ an argument, since no config field holds it); \code{per\_mode\_spectrum(lam, sx2, sy2, sig2, r, delta) -> (c, d)} (the port of \code{spectrum\_d}, on the physical branch); \code{predicted\_spectra(cfg, dataset) -> Prediction(d (K, N), null\_floor, c (K,), tau (K,), params)}; \code{predicted\_summary(d) -> (q\_half, share0, lag1)}, which may read $q_{1/2}$ because a prediction is a mean spectrum; and \code{ansatz\_free\_kernel(cfg, dataset) -> C\_tilde (P, P)} from the master equation, for any teacher. \code{theory\_parameters} and \code{predicted\_spectra} refuse a non-orthonormal teacher; \code{network\_parameters}, which every function taking a \code{Config} calls (\code{ansatz\_free\_kernel} included), refuses a float32 run; \code{per\_mode\_spectrum} takes arrays and checks neither. Tests: the module reproduces the limits note's quoted numbers for predictions 1--3 ($\tau_1 = 0.00814$, shares $0.157/0.156$, $0.077/0.111$, $0.343/0.164$; the ratio $10.08$ at $q = 0$ of P1 and $53.5$ of P3) from a \code{Config} built as in Section~\ref{sec:experiments}. The mean-field solver itself is not changed; a spatial extension of it is a separate, later item.

\emph{Limitation: the theory is the limit $N\to\infty$ at fixed $n$.} An exact Gibbs sampler for the linear network agrees with SGLD and shows what \code{predicted\_spectra} omits. Integrating the readout out exactly leaves $\frac n2\log\det\tilde C + \frac N{2k_1}\tr(Y^\top\tilde C^{-1}Y)$ in the action; the mean-field per-$q$ equations keep the second term and drop the first, which, linearised, adds $+\frac nN\lambda_q\tilde C^{-1}$ to each per-$q$ precision: an $O(n/N)$ correction. Its size, as the ratio of the exact to the predicted $q = 0$ power at $\beta = 10$ on the tiny geometry: $0.53$, $0.69$, $0.81$, $0.90$, $0.94$ at $N = 16$, $32$, $64$, $128$, $256$, growing with $n$; on the full preset ($N = 64$, $n = 8$) at $\beta = 10$, $0.57$; at $\beta = 0.25$ the theory gets the sign of the data effect wrong (by about $5\%$). The notes' theory is this $N\to\infty$ limit at fixed $n$. No code is added for the correction: the experiments are run in the regime $N\gg n$ where the theory is accurate (Section~\ref{sec:protocol}(vi); W10, W11), and the predicted numbers of this plan, computed at $N = 64$, are recomputed with \code{predicted\_spectra} at the $N$ actually run.

\emph{Cache.} None (pure, no artefacts).

\subsection{D5: names and ranges}
\label{sec:D5}
(a) \code{net.g} is a hidden-layer gain, unrelated to the notes' rate coefficient $r$ (which the dictionary maps to \code{rate\_lambda} $= r$); no code change is needed for $L = 1$, and the preset pins \code{g = 1.0}. (b) $k_0$, $k_1$ are kernel amplitudes with \code{length\_scale = 0} and \code{weight\_lambda = 1}; a common scale $k_0 = k_1 = k$ is \code{weight\_lambda} $= 1/k$ at unit amplitudes (Section~\ref{sec:dictionary}), which is what prediction 8 uses. (c) $\ell$ is in units of \code{box}; the preset sets \code{box} $= N$ so that \code{length\_scale} is in sites. (d) $n = m$ is forced by \code{linear\_teacher}. (e) $K$ latents are \code{data.n\_latents}; the number of features $m$ is \code{data.D\_in} (in \code{ExpConfig}: \code{linear\_teacher.n\_features}, \path{fct/exp.py:153}).

\subsection{Not discrepancies, but protocol}
\label{sec:protocol}
(i) \emph{Finite step size.} The stationary variance under SGLD is $1/(\beta c(1 - \mathrm{lr}\,c/2))$ on a quadratic of curvature $c$ (Section~\ref{sec:loss}), and an Euler step diverges once $\mathrm{lr}\,c\ge2$. Under the whitened parametrisation of D1 every prior curvature is $\lambda_w/\beta$. The data term adds, along the directions the data spans, a curvature that does \emph{not} shrink with $\beta$ and grows with the weight scales: $\xi^{(0)}$ feels $\|V\|^2\propto k_1$ and $\xi^{(1)}$ feels $\|h\|^2\propto k_0s_1^2$, so with the common scale $k = 1/\lambda_w$ at unit amplitudes
\begin{equation}
  c_{\mathrm{data}}\approx C\,\frac{\lambda_0s_1^2}{N\,m\,n\,\lambda_w},\qquad C = O(1),
  \label{eq:cdata}
\end{equation}
($\lambda_0s_1^2/(Nmn) = 0.52$ at the full size, $0.37$ on the smoke ring). Measured on exact Hessians of $L$ in whitened coordinates at prior samples (orthonormal teacher, $Z\sim\mathcal U[0,1]$, three data seeds; the review of this note), $\lambda_{\max}(\nabla^2L) = 4.8$ (full P1, $\beta = 0.25$: the prior dominates), $0.8$ (P3, $\beta = 10$), $5.1$ (P8, $k = 10$), $50$ ($k = 100$), $1.85$ (smoke S3), $4.5$ (S8, $k = 10$), $45$ (S8, $k = 100$): $C\approx1.5$ at the full size and up to $5$ on the smoke ring. The rule used throughout, for $\mathrm{lr}\,\lambda_{\max}\le0.1$ (at most $5\%$ inflation of the stiffest mode, $1\%$ of the prior ones):
\begin{equation}
  \mathrm{lr} = \min\Big(\frac{0.02\,\beta}{\max(1,\lambda_w)},\ \frac{0.1}{\lambda_w/\beta + 1/\beta + \hat c_{\mathrm{data}}}\Big),\qquad
  \hat c_{\mathrm{data}} = \frac{C\lambda_0s_1^2}{Nmn\lambda_w},\ C = 2\ (\text{full}),\ 5\ (\text{smoke}),
  \label{eq:lr-rule}
\end{equation}
where the term $1/\beta$ is a margin (no parameter has that curvature at $\sigma = 0$; keeping it leaves every \code{lr} below as computed), and the whitened path keeps a loud check of $\mathrm{lr}\,\lambda_{\max}$ at run start (D1). The rule is for the whitened path; the penalty path's step is set by its own check (D1 and (v) below). At $\beta = 0.25$ the prior term sets the step ($\mathrm{lr} = 0.005$); at $\beta = 10$ the data term does: $0.08$ at $k = 1$, $0.01$ at $k = 10$, $0.001$ at $k = 100$, not the $0.2$ a prior-only rule gives, at which $k = 100$ diverges at the first step. A data-spanned direction then relaxes in $1/(\mathrm{lr}\,c_{\mathrm{data}})\approx10$ steps; a prior-only direction in $\beta/(\lambda_w\mathrm{lr})$ steps, up to $10^6$ at $k = 100$, which is why $\xi$ is initialised from its prior (D1): those directions start stationary and need no relaxation, and the slowest remaining directions are the weakly data-driven ones (a contrast mode at a $q$ with $\lambda_q/\lambda_0 = 0.1$ has $c\approx0.1c_{\mathrm{data}}s_2^2/s_1^2$, i.e.\ $\sim2\times10^3$ steps at $k = 100$). (ii) \emph{Sampling: one sample per member.} One SGLD chain is not trusted to visit independent networks, so each of the $B$ members of an ensemble contributes exactly one sample, its final state, and $B$ is the sample size. Every predicted statistic, the $q = 0$ share $d_0/\sum_qd_q$, $R(\Delta)/R(0)$, the Gaussian width, the ratio $\hat d^{(1)}_q/\hat d^{(2)}_q$ and $q_{1/2}$, is a function of the mean spectrum $d^{(k)}_q$, and a ratio of means is not a mean of ratios: the mean over members of per-member shares or lags is biased (simulated from the theory's law: P1's share $0.145$ against $0.157$, $R(6)/R(0)$ low by $0.05$, P3's share at $\beta = 10$ $0.278$ against $0.343$, S1's share $0.273$ and $R(1)/R(0)$ $0.859$ against $0.314$ and $0.883$), and $q_{1/2}$ of one draw is noise, $d_0$ and $d_{N/2}$ being single $\chi^2_1$ draws (\code{half\_power\_q} also raises on a rising spectrum). So each of them is read off the final-state spectrum averaged across the $B$ members, one value per setting (at $K = 3$ rows 1 and 2, both contrasts, may be pooled for mode 2), with an error bar from a bootstrap over members: resample the $B$ members with replacement, average, recompute. Per-member values serve only as spreads and diagnostics; the hexatic moment of P6 is itself a mean, of $e^{6i\theta}$, and its complex average across members is its estimator. More samples come from \code{fct exp extend ROOT --by K}, which trains only the new members of every setting and appends them (member $b$'s initial weights do not depend on $B$; its SGLD noise comes from the ensemble's single global stream, so extended members equal those of a from-scratch run in distribution, not bit for bit); for that \code{B} is a shared field of the experiment, never an axis. The records along a chain are a diagnostic of burn-in and mixing only: \code{average\_records}, and \code{record\_sem} and \code{integrated\_autocorr\_time} (Geyer's initial monotone sequence), which need at least $200$ records after burn-in and return a \code{reliable} mask, false wherever $n < 20\tau$, $\tau\le0$ or Geyer's sequence does not end before $n/2$; compared across burn-in values and with the spread across members they show whether the chains have equilibrated, and they never enter an estimator.

\emph{What $B$ buys.} Under the theory's law the mean over $B$ members of $|\tilde t_k(q)|^2$ is $d_qG_q$ with $G_q\sim\mathrm{Gamma}(B, 1/B)$ for $0 < |q| < N/2$ and $\mathrm{Gamma}(B/2, 2/B)$ at $q = 0, N/2$ (one member's draw is $d_q$ times an $\mathrm{Exp}(1)$, or a $\chi^2_1$ at $q = 0, N/2$); the maps are real, so $\tilde t_k(-q) = \overline{\tilde t_k(q)}$ and folding $\pm q$ averages equal numbers. Simulated from the predicted spectra at $\sigma = 0$ (Appendix~\ref{app:numbers}), the standard errors of the member-averaged estimators at $B = 16$, all falling as $B^{-1/2}$ and with biases below $0.01$, are: the share $0.049$ (P1, each mode), $0.085$ and $0.051$ (P3 at $\beta = 10$), $0.083$ (smoke S1), and $0.068$ for the difference of P1's two shares; $R(\Delta)/R(0)$ in P1 $0.003$, $0.012$, $0.025$, $0.039$, $0.052$, $0.063$ at $\Delta = 1,\dots,6$, and $R(1)/R(0)$ $0.017$ in S1; the width at $\Delta = \ell$ $0.09$, $0.26$, $0.73$ sites at $\ell = 2, 4, 8$; the ratio of the two modes $0.62$ relative at $q = 0$ and $0.39$ at $|q|\ge1$. $q_{1/2}$ of P1's mode 1 ranges over $0$--$5$ at $B = 16$ (predicted $2$--$3$), $1$--$4$ at $64$ and $2$--$3$ at $256$; P2's mode 1 gives its predicted $6$ in $44\%$, $75\%$, $98\%$ of draws at $B = 16$, $64$, $256$. A criterion with margin $\epsilon$ is a three-standard-error test when the standard error is at most $\epsilon/3$, i.e.\ $B\ge16\,(3s_{16}/\epsilon)^2$ with $s_{16}$ the standard error at $B = 16$. So $R(1)/R(0)$ to $0.05$ and the widths of P7 within $20\%$ hold at $B = 16$ ($B\approx30$ at $\ell = 8$); P3's $q = 0$ ratio, predicted $54$, falls below $20$ in $2.9\%$ of draws at $B = 16$ and $10^{-4}$ at $B = 64$; $R(\Delta)/R(0)$ to $0.05$ out to $\Delta = 6$ needs $B\approx230$; a difference of shares to $0.03$ needs $B\approx750$; the ratio within $20\%$ at one $|q|$ needs $B\approx550$ ($1400$ at $q = 0$). The band share of $|q|\le3$, $\sum_{|q|\le3}d_q/\sum_qd_q$, has a standard error of $0.029$ at $B = 16$ in P1 ($0.041$ for the difference of the two modes, $0.014$--$0.045$ elsewhere), and $R(1)/R(0)$ $0.002$--$0.006$. So the falsifiers of Section~\ref{sec:experiments} are ring-averaged statistics, $R(\Delta)/R(0)$, the band share and the widths, which need $B\approx16$--$230$; single-$q$ shares, single-$q$ ratios and $q_{1/2}$ are secondary criteria, for which $B$ is raised with \code{exp extend} when they are wanted. $B$ is set per falsifier by these numbers, with $B\ge16$ the minimum (P6's); the presets item sets the concrete values. (iii) \emph{Finite $N$.} The predictions are evaluated on the $N = 64$ ring exactly (no continuum approximation), and finite-size corrections to the theory are $O(e^{-N/\xi})$ (\path{continuum_limit.tex}, Section~3.3); the mean-field step itself, the self-averaging of $C$ over the posterior, is $O(N^{-1/2})$ and is the reason to confirm prediction 3 at $N = 256$ (Section~\ref{sec:P3}). (iv) \emph{The $q = 0$ share is $N$-specific} (Summary): at fixed $\tau_1 = 0.835$ it is $0.343, 0.172, 0.086, 0.043$ for $N = 64, 128, 256, 512$, while $R(1)/R(0) = 0.985$ does not move.

(v) \emph{The two parametrisations as an equivalence test.} The parametrisation of the spatial prior, \code{"penalty"} or \code{"whitened"} (\code{spatial.parametrisation}), is a factor of every run of W10 and W11: the two paths sample the same stationary law up to the eigenvalue floor (D1), so their final-state statistics must agree, and a disagreement is a bug in one of them or a chain that has not reached stationarity. Their dynamics differ (D1, as implemented), so they are not run at one \code{lr}: each parametrisation runs at the largest \code{lr} its own run-start check accepts (for the whitened path \eqref{eq:lr-rule}, under its Hessian check; for the penalty path just below $2\beta f\lambda_0/\lambda_w$ by its prior check, $5\times10^{-6}$ at $\beta = 0.25$ and $2\times10^{-4}$ at $\beta = 10$ with $f = 10^{-6}$, $\lambda_w = 1$) and for as many steps as it needs to reach stationarity, judged by the burn-in diagnostics of (ii). A known property of the penalty path: its check refuses only at $\mathrm{lr}\,c\ge2$ for its stiffest prior curvature $c$, the Euler divergence edge, where the whitened check refuses at $\mathrm{lr}\,\lambda_{\max}\ge1$; at a step inside that edge the floored modes' stationary variance is inflated by $1/(1 - \mathrm{lr}\,c/2)$, without bound as the edge is approached; and the penalty path has no data-curvature check. The presets therefore take a stated fraction just inside each path's bound, and the inflation that fraction implies for the floored modes is part of the comparison. Both start from their prior (\code{net.init = "prior"}), so only the data-driven directions relax, in $\sim1/(\mathrm{lr}\,c)$ steps: $10$ to a few $10^3$ on the whitened path (Section~\ref{sec:protocol}(i)), and on the penalty path at least the ratio of the two steps more ($10^3$ at $\beta = 0.25$), so $\sim10^4$--$10^7$ steps; the step counts are therefore set per path. The floor $f$ = \code{min\_rel\_eig} separates the two posteriors by the shifts of D1 (at $f = 10^{-3}$ P2's mode-1 $R(1)/R(0)$ falls from $0.922$ to $0.812$; at $10^{-4}$ by $0.012$, with P1 and P3 within $0.001$; at $10^{-6}$ nothing moves by $10^{-3}$, but the penalty step must then be below $5\times10^{-6}$ at $\beta = 0.25$); the real runs use $f = 10^{-6}$ (D1). \code{theory\_parameters} gives each path the $\lambda_q$ it trains with, so each path is compared with its own prediction and the two with each other up to the predicted shift. The comparison is of the final-state statistics of (ii) across the $B$ members of each path, as a difference with its standard error. Each experiment is realised either as two ROOTs, one per parametrisation, each with its own \code{--set lr=} and \code{--set n\_epochs=} (the simplest, since both differ), or as one ROOT with \code{--axis 'spatial.parametrisation=["penalty","whitened"]'} paired with a \code{train.lr} axis of which only the diagonal is run, as P3 pairs $\beta$ with \code{lr}; the second shares \code{n\_epochs}, which must then be the penalty path's.

(vi) \emph{The large-$N$ regime.} The per-mode theory is the mean-field limit $N\to\infty$ at fixed $n$, and at finite $N$ it misses an $O(n/N)$ correction (D4): on the full preset ($N = 64$, $n = 8$) at $\beta = 10$ the exact $q = 0$ power is $0.57$ of the predicted. Experiments are therefore compared with the theory only where $N\gg n$ and the theory is accurate. The regime of each preset ($N$, $n$, $\beta$, $r$, $k$) is chosen once, by the \code{linear-ring-presets} item, so that the mean-field theory is accurate there; that item checks it off-repo against an exact Gibbs computation for the linear network and records the exact-versus-predicted numbers in the preset's description. No sampler is run per experiment, and \code{fct} gains no code for it. A run outside a documented regime is not compared against \code{predicted\_spectra}. At $N = 64$, $n = 8$ the theory fails at both $\beta$ checked (at $\beta = 10$ the $q = 0$ power is $0.57$ of the predicted; at $\beta = 0.25$ the data effect has the wrong sign), so the presets item raises $N$ (with \code{spatial.box} $= N$, so $\ell$ stays in sites) until it holds.

\section{Dictionary: theory symbol to \texorpdfstring{\code{Config}}{Config} field}
\label{sec:dictionary}
\begin{table}[H]
\centering\footnotesize
\begin{tabular}{@{}>{\raggedright\arraybackslash}p{2.1cm}>{\raggedright\arraybackslash}p{2.5cm}>{\raggedright\arraybackslash}p{7.0cm}>{\raggedright\arraybackslash}p{2.9cm}@{}}
\toprule
Symbol & Meaning & \code{Config} field (value for the common setting) & Reference\\
\midrule
$N$ & hidden units on the ring & \code{net.N} (64) & \path{model.py:46}\\
$L$ & hidden layers & \code{net.L} (1) & \path{model.py:47}\\
$\phi$ & activation & \code{net.phi} (\code{"id"}) & \path{activations.py:21}\\
$m$ & inputs & \code{data.D\_in} (8); \code{ExpConfig.linear\_teacher.n\_features} & \path{tasks.py:88}\\
$n$ & outputs & $= m$ for \code{linear\_teacher} (8) & \path{tasks.py:264}\\
$P$ & patterns & \code{data.P} (256) & \path{config.py:49}\\
$K$ & latents & \code{data.n\_latents} (3) & \path{config.py:60}\\
$[0,a]$ & latent interval & \code{data.latent\_min} (0.0), \code{data.latent\_max} (1.0) (W1) & \code{DataConfig}\\
$W_X$, $W_Y$ orthonormal & teacher & \code{data.teacher = "orthonormal"} (W2) & \\
$\ell$ (sites) & kernel width & \code{spatial.hidden.kernel.length\_scale} (4.0) with \code{spatial.box} $= N$ & \path{spatial.py:201--221}\\
$K_{i,i} = 1$ & hidden amplitude & \code{spatial.hidden.kernel.amplitude} (1.0) & \path{spatial.py:221}\\
$k_0$ & input prior scale & \code{spatial.input.kernel.amplitude} (1.0), \code{length\_scale} 0 & \path{spatial.py:218--219}\\
$k_1$ & readout prior scale & \code{spatial.output.kernel.amplitude} (1.0), \code{length\_scale} 0 & \\
$k_0 = k_1 = k$ & common scale & \code{train.weight\_lambda} $= 1/k$ at unit amplitudes & \path{train.py:45--54}\\
$\lambda_w$ otherwise & & \code{train.weight\_lambda} (1.0) & \\
periodic grid & embedding & \code{spatial.*.embedding.kind = "grid"}, \code{d = 1} & \path{spatial.py:160}\\
exact $\lambda_q$ & no floor & \code{spatial.parametrisation} $\in\{$\code{"whitened"}, \code{"penalty"}$\}$, both run (protocol (v)); \code{min\_rel\_eig} $10^{-6}$ & D1\\
initial law & the prior & \code{net.init = "prior"} (shared field \code{init}) & D1\\
$\beta$ & inverse temperature & \code{train.beta} (per experiment); \code{train.optimizer = "sgld"} & \path{sgld.py:60}\\
$\kappa$, $\delta$ & $\beta/(Pn)$, $Pn/(\beta k_1)$ & derived & \\
$r$ & rate penalty & \code{train.rate\_lambda} $= r$ from \code{LOSS\_VERSION = 5} (Section~\ref{sec:loss}) & \path{train.py:93}, the unscaled term\\
$\sigma^2$ & activity noise & none: $\sigma^2 = 0$ (D2) & \\
gain & $1/\sqrt N$ only & \code{net.g} (1.0) & \path{model.py:92}\\
SGLD step & $\mathrm dt$ & \code{train.lr} by \eqref{eq:lr-rule} (whitened); penalty by its check & Section~\ref{sec:protocol}\\
steps & & \code{train.n\_epochs}; record every \code{train.monitor\_int} with \code{train.monitor\_kind = "spectrum"} (W7) & \path{train.py:427}\\
sample size $B$ & members & \code{net.B}, the shared field \code{B} ($\ge16$, set per falsifier), never an axis; one sample per member, its final state; more with \code{exp extend} & \path{model.py:103}\\
\bottomrule
\end{tabular}
\caption{The common setting of the limits note (Section~4 there) in \code{Config} fields. Everything not listed keeps the library default; \code{dtype = "float64"}, \code{device} as available.}
\label{tab:dictionary}
\end{table}

Derived for the common setting: $\lambda_0 = \sqrt{2\pi}\,\ell = 10.03$, $\bar\lambda = 1$, $\bar\lambda_2 = \sqrt\pi\,\ell = 7.09$; $s_1^2 = 256\cdot10/12 = 213.3$, $s_{2,3}^2 = 21.33$; $\hat s_{X,1}^2 = 26.67$, $\hat s_{X,2}^2 = 2.667$; $\hat s_{Y,k}^2 = s_k^2$ at $k_1 = 1$; $\delta = 2048/\beta$.

\section{Measurement: functions, inputs, outputs, maths}
\label{sec:measure}
The theory predicts (limits note, Section~3.1) that the tuning map $t_{i,k} = u_k^\top h_i$, $h_i\in\R^P$ the activity of neuron $i$ over patterns and $u_k$ the $k$-th left singular vector of $Z$, is a stationary Gaussian field on the ring with power spectrum $d^{(k)}_q$ in the unitary DFT, $\langle|\tilde t_k(q)|^2\rangle = d^{(k)}_q$, independent across $k$, and that the $P - K$ null directions have the flat spectrum $\sigma_0^2 = \sigma^2/(1 + r\sigma^2)$, which is $0$ at $\sigma = 0$ (D2): the null row then checks the projection, not the theory. The expectation is over the posterior; the estimators below draw one sample of it per ensemble member (Section~\ref{sec:protocol}(ii)).

\paragraph{Pure analysis (\code{fct/spectrum.py}, W6, done).} Pure functions on tensors, any floating dtype and device, leading batch dimensions (members, records) passed through; validated loudly (shapes, \code{K < P}, finite inputs).
\begin{itemize}[nosep]
  \item \code{latent\_basis(factors) -> LatentBasis(U (P, K), s2 (K,), Z (P, K))}: $Z$ stacks \code{z0..z\{K-1\}} by index; thin SVD $Z = USV^\top$ with singular values descending, so column 0 is the mean direction ($u_1$) and columns $1..K-1$ the contrasts; \code{s2} $= S^2$ (the $s_k^2$ the theory needs, from the data, so finite-$P$ fluctuations are included). Sign convention: $\mathbf 1^\top u_k\ge0$. Raises unless $u_1$ is the mean direction (a gap $s_1^2/s_2^2\ge2$, against $3K + 1$ for latents on $[0, a]$, and a cosine of $u_1$ with $\mathbf 1$ at least half its $[0, a]$ value): symmetric latents have a degenerate spectrum and a random SVD frame, and are refused.
  \item \code{tuning\_maps(h (..., P, N), U) -> (..., K, N)}: $t = U^\top h$; \code{power\_spectrum(t)}: $|\tilde t_k(q)|^2$ with $\tilde t = \code{fft}(t)/\sqrt N$ (unitary; $q$ in FFT order, $0,1,\dots,N/2,\dots,-1$; no mean subtracted); \code{null\_power(h, U) -> (..., N)}: the mean spectrum of the $P - K$ directions orthogonal to $\col(Z)$; \code{mode\_spectra(h, U) -> (..., K+1, N)}: the $K$ latent rows and the null row, as the monitor records them.
  \item \code{fold\_q(d) -> (..., N/2 + 1)}: average of $q$ and $-q$, $|q| = 0,\dots,N/2$ (for one real map the two are equal).
  \item \code{autocorrelation(d)}: $R(\Delta) = \frac1N\sum_qd_qe^{2\pi iq\Delta/N}$ (\code{ifft(d).real}).
  \item \code{spectrum\_summary(d) -> Summary(share0, lag1, d0, d\_nyq, R)}: \code{share0} $= d_0/\sum_qd_q$, \code{lag1} $= R(1)/R(0)$, $d_0$, $d_{N/2}$ and $R$; defined on any spectrum; for the predictions it is applied to the member-averaged final-state spectrum (protocol (ii)), per member only for spreads and diagnostics.
  \item \code{half\_power\_q(d)}: $q_{1/2}$, the largest $|q|$ with $d_q - d_{N/2}\ge\frac12(d_0 - d_{N/2})$ (the limits note's definition, with $d_{N/2}$ as the floor); a statistic of a mean spectrum only, and it raises where $d_0 < d_{N/2}$. It is applied to the final-state spectrum averaged across the $B$ members, one value per setting, with a bootstrap over members for the error bar.
  \item \code{gaussian\_width(R, ell) -> (...)}: $\sqrt{-\ell^2/(2\ln(R(\ell)/R(0)))}$ for an integer lag $\ell$ (prediction 7); it raises for the whole batch if any entry has $R(\ell)/R(0)$ outside $(0, 1)$, so it is applied to the member-averaged $R$.
  \item \code{hexatic\_moment(t2, t3) -> complex (...)}: $m = \langle e^{6i\theta_i}\rangle$ over neurons, $\theta_i = \operatorname{atan2}(t_{i,3}, t_{i,2})$ in the contrast plane $(u_2, u_3)$; the statistic of prediction 6 is $|m|$, and the complex $m$ is what is averaged across members (the modulus of a mean is not the mean of moduli). \code{theta\_histogram(t2, t3) -> (..., 36)}: counts on fixed bins of $(-\pi, \pi]$, summable across members. \code{contrast\_angle\_statistic} ($\langle\cos4\theta\rangle$) also exists but is not used (Section~\ref{sec:P6}).
  \item \code{member\_mean\_sem(values (B, ...))}: mean and $\mathrm{sd}/\sqrt B$ across members, for the spread of per-member values (a diagnostic) and for P6's complex hexatic moment; the predicted statistics, being ratios of the mean spectrum, take a bootstrap over members instead (Section~\ref{sec:protocol}(ii)).
  \item Diagnostics of one chain: \code{average\_records(records, burn\_in)}; \code{integrated\_autocorr\_time(records, burn\_in) -> (tau, reliable)} (Geyer's initial monotone sequence) and \code{record\_sem(records, burn\_in) -> (sem, tau, reliable)}. Both need at least $200$ records after burn-in (fewer raise) and mark an element unreliable, its value NaN, where $n < 20\tau$, $\tau\le0$ or the sequence does not end before $n/2$; callers read \code{reliable}. None of them is an estimator (Section~\ref{sec:protocol}(ii)).
\end{itemize}

\paragraph{The monitor (\code{fct/monitor.py}, W7, done).} \code{make\_spectrum\_monitor(dataset, cfg) -> Monitor}, selected by \code{train.monitor\_kind = "spectrum"} (hashed only as \code{"spectrum"}; in an experiment the shared field \code{monitor\_kind}, which cannot be an axis). It builds \code{latent\_basis(dataset.factors)} once, so data without a mean direction fail before training, and an experiment with the spectrum monitor is validated at load: every setting must be \code{linear\_teacher} with latents \code{latent\_basis} accepts. Every \code{monitor\_int} steps it runs \code{net(dataset.x)} under no-grad (there is no input noise), takes the pre-activations of every hidden layer (for \code{"id"} equal to the rates) and records, per layer, \code{spectrum} ($(K+1, N)$ float32, \code{mode\_spectra}) and, for $K\ge3$, \code{hex} (the real and imaginary parts of \code{hexatic\_moment} of rows 1 and 2) and \code{theta\_hist} (36 counts); the history keys are \code{monitor\_spectrum}, \code{monitor\_hex}, \code{monitor\_theta\_hist}, with a member axis for an ensemble. A call draws from no RNG, so monitoring leaves the training stream unchanged. The records serve the burn-in diagnostics; the estimators average the members' final states, i.e.\ \code{mode\_spectra} of the final network's pre-activations on the training set (the last record, \code{monitor\_int} steps earlier, is an equally valid sample).

\paragraph{Theory side (\code{fct/linear\_theory.py}, W5, done).} \code{predicted\_spectra(cfg, dataset)} returns $d^{(k)}_q$ in the same FFT order and the null floor ($0$ at the default $\sigma^2 = 0$), computed with the \emph{measured} $s_k^2$ of the dataset's latents (the limits note's formulas are the $P\to\infty$ values; at $P = 256$ the fluctuations are $O(P^{-1/2})\approx6\%$ on $s_k^2$, which enters $\sX$, $\sY$ linearly) and with the $\lambda_q$ of the parametrisation the run trains with (D4). The note's predicted numbers below use the formula values; the implementing session should quote both.

\paragraph{Plotting (\code{fct/draw.py}, W9, to do; new \code{plot\_*} taking an \code{ax}).} \code{plot\_spectrum(d\_measured (K+1, N/2+1), sem, d\_theory, ax)} (log $y$ versus $|q|$, one colour per mode, theory as lines, measurement as points with error bars, the null row dashed); \code{plot\_spectrum\_ratio}; \code{plot\_autocorrelation(R, kernel\_profile, ax)}; \code{plot\_share\_sweep(values, shares (n, 2), theory, ax)}; \code{plot\_angle\_histogram(theta, ax)}. A notebook builder \path{notebooks/build_linear_ring.py} in the style of \path{build_spatial_exp.py} (parameters \code{ROOT}, \code{SETTING}, \code{BURN\_IN}) shows, per setting: loss curves, the burn-in diagnostics of the records, the member-averaged final-state spectra against theory, the ratio, the autocorrelation, the \code{"penalty"} and \code{"whitened"} settings side by side, and prints the summary table (every statistic of the member-averaged final-state spectrum, $q_{1/2}$ included, with its bootstrap error over members) against the predicted numbers.

\section{The experiments}
\label{sec:experiments}

\subsection{Presets and invocation}
\label{sec:presets}
Two new presets (JSON files in \path{fct/presets/}, complete \code{ExpConfig}s; \path{docs/design.md} ``Presets''), defined by W8 (\code{linear-ring-presets}, in progress); the values below are the plan's, and W8's files are authoritative where they differ. Fields of \code{ExpConfig} (\path{fct/exp.py:173--232}): the shared network and loss fields reach every setting (\code{\_base\_config}, 478--501), \code{--set} overrides any shared entry by dotted key and \code{--axis} adds an axis over any \code{Config} field (\code{cmd\_exp\_init}, \path{fct/cli.py:384--413}); the fields the experiment owns (\code{EXP\_OWNED\_AXES}, \path{fct/exp.py:106--117}: \code{data.P}, seeds, log/save/monitor intervals, and on \code{sigma0} \code{train.monitor\_kind}) cannot be axes. On \code{sigma0} the fields of W1--W3 and W7 exist and are set in the presets: \code{data.latent\_min}, \code{data.latent\_max} and \code{data.teacher} through the \code{linear\_teacher} parameters, \code{spatial.parametrisation}, and the shared fields \code{init} and \code{monitor\_kind}; \code{B} is a shared field (never an axis, so that \code{exp extend} can grow it).

\code{linear\_ring} (full):
\begin{verbatim}
N 64, P 256, n_epochs 200000, lr 0.005, log_int 1000, save_int 10000,
monitor_int 50, monitor_kind "spectrum",
monitor_bins 10, monitor_shuffles 5   (required by ExpConfig, unused here),
linear_teacher {n_features 8, n_latents 3, latent_min 0.0, latent_max 1.0,
                teacher "orthonormal"},
object_action {grid_size 5, max_objects 3}, what_where {n_motifs 8}   (required, unused)
seeds {data 0, net 0, train 0}, dtype float64, device auto, task linear_teacher,
L 1, g 1.0, B 16, init "prior", phi "id", weight_lambda 1.0, rate_lambda 0.0,
optimizer "sgld", beta 0.25,
spatial {d 1, box 64.0, parametrisation "whitened", min_rel_eig 1e-6,
         input  {embedding grid, kernel {length_scale 0.0, amplitude 1.0}},
         hidden {embedding grid, kernel {length_scale 4.0, amplitude 1.0}},
         output {embedding grid, kernel {length_scale 0.0, amplitude 1.0}}},
axes {"train.beta": [0.25]}   (one-value axis; experiments replace it)
\end{verbatim}
\code{B 16} is the smallest ensemble prediction 6 accepts ($64B\ge10^3$ neuron-samples); the other predictions need more members for their criteria (Section~\ref{sec:protocol}(ii)), added with \code{exp extend}. \code{B} is set per falsifier by Section~\ref{sec:protocol}(ii) ($B\ge16$); \code{min\_rel\_eig 1e-6} matters on the penalty path only (D1). \code{N 64} is the size the predicted numbers are computed at; the runs need $N\gg n$ (Section~\ref{sec:protocol}(vi)), and the presets item sets the $N$ that satisfies it. \code{linear\_ring\_tiny} (smoke): the same with \code{N 16, P 32, n\_features 4, n\_latents 2, n\_epochs 10000, lr 0.005, log\_int 100, save\_int 1000, monitor\_int 10, B 16, device cpu}, optionally \code{spatial.min\_rel\_eig 1e-3} (a pipeline check against its own floored prediction), \code{spatial.box 16.0} (so that \code{length\_scale} stays in sites: with the inherited $64$ the spacing would be $4$ and $\ell = 2$ half a site) and hidden \code{length\_scale 2.0}. Its predicted numbers are in Section~\ref{sec:smoke}. The description of each preset states the dictionary of Table~\ref{tab:dictionary} and the step rule \eqref{eq:lr-rule}: \emph{every value of $\beta$ or $\lambda_w$ must come with the matching \code{lr}}, and the run-start check of D1 refuses a step that is too large. $\beta$ is always given as the axis \code{--axis 'train.beta=[...]'}: an experiment without axes is rejected (\code{\_validate\_axes}, \path{fct/exp.py:555--556}), \code{axes = None} would mean the modularity axes (\code{axes\_of}, \code{:456}), and an axis wins over the shared field, so \code{--set beta=} would be ignored. A single value is \code{--axis 'train.beta=[10]'}; \code{lr} goes with \code{--set lr=} (one value) or a parallel \code{--axis 'train.lr=[...]'} of which only the diagonal is run.

\emph{Both parametrisations.} Every invocation below is written for the whitened path. Its penalty twin is the same command in a sibling ROOT (\code{ROOT/p1-penalty} beside \code{ROOT/p1}) with \code{--set spatial.parametrisation=penalty}, \code{--set lr=} at a stated fraction just inside the penalty path's bound $2\beta f\lambda_0/\lambda_w$ (Section~\ref{sec:protocol}(v)) and \code{--set n\_epochs=} long enough for stationarity; the two are compared on their final-state statistics.

Burn-in: the first $10^4$ steps of the full version and the first $2\times10^3$ of the smoke (200 records each) are discarded by the diagnostics (\code{BURN\_IN}); with $\xi$ initialised from its prior only the data-driven directions have to relax, in $10$ to $2\times10^3$ steps on the whitened path (Section~\ref{sec:protocol}(i)), and the burn-in of the penalty path scales with its step count. Run time: of order a millisecond per step at the full size on a few CPU threads (Appendix~\ref{app:numbers}); with the monitor, of order $10$--$30$ minutes per whitened setting on a CPU at $B = 4$ and up to four times that at $B = 16$, under a minute for the smoke. SLURM: \code{bash jobs/fct\_exp.sh ROOT --submit}; after \code{fct exp extend ROOT --by K} the same command trains only the new members.

Seeds: \code{seeds.data = 0} fixes $Z$, $W_X$, $W_Y$ (the predicted $s_k^2$ are then fixed numbers, reported by \code{latent\_basis}); \code{seeds.net = 0} the initial weights (member $b$ from its own seed, independent of $B$); \code{seeds.train = 0} the SGLD noise. Error bars come from the $B$ members (independent chains on the same data, one sample each, \path{docs/design.md} ``Ensembles''); a second data seed (\code{--set seeds.data=1}) in a fresh ROOT checks that nothing depends on the particular $Z$.

\subsection{The smoke version}
\label{sec:smoke}
Ring $N = 16$, $\ell = 2$ ($\lambda_0 = 5.013$, $\lambda_q/\lambda_0 = 1, 0.735, 0.291, 0.062, 0.0072, \dots$ at $|q| = 0,1,2,3,4$), $P = 32$, $m = n = 4$, $K = 2$, $a = 1$: $s_1^2 = 32\cdot7/12 = 18.67$, $s_2^2 = 2.667$, $\hat s_{X,k}^2 = 4.667, 0.667$, $\delta = 128/\beta$. Predicted at $\sigma^2 = 0$ (computed as in Appendix~\ref{app:numbers}; $|q| = 0,\dots,8$):
\begin{itemize}[nosep]
  \item S1 ($r = 0$, $\beta = 0.25$): $\tau_1 = 0.0078$, $\tau_2 = 1.7\times10^{-4}$; $q_{1/2} = 1$ both; shares $0.314$, $0.313$; $R(1)/R(0) = 0.883$ both; ratio $7.05$, $7.03$, $7.00$ and exactly $7$ beyond (the ratio is $\hat s_{X,1}^2/\hat s_{X,2}^2 = 7$ for $K = 2$; with no floor both modes share the $\lambda_q$ tail). Falsified if the shares differ by more than $0.03$ or the ratio over $|q|\le3$ is outside $7\pm1$.
  \item S2 ($r = 0.5$, else as S1; $\gamma_1 = 11.7$, $\gamma_2 = 1.67$): $q_{1/2} = 2$ both (the ring is too short to separate them), but the shares $0.175$ (mode 1) versus $0.241$ (mode 2) and $R(1)/R(0) = 0.757$ versus $0.839$ show the whitening; the ratio rises monotonically, $1.47$, $1.63$, $2.36$, $4.47$, $6.53$, $6.97$, $7.00$, $7.00$, $7.00$ (no interior maximum, as in P2 at $\sigma = 0$). Falsified if mode 1's share is not below mode 2's by at least $0.03$, or $\hat d_0^{(1)}/\hat d_0^{(2)} > 3$. ($r = 0.1$ gives shares $0.229/0.289$ and $\hat d_0^{(1)}/\hat d_0^{(2)} = 2.80$, also usable.)
  \item S3 ($\beta = 10$, $r = 0$; $\mathrm{lr} = 0.05$ by \eqref{eq:lr-rule} with $\hat c_{\mathrm{data}} = 1.8$): $\tau_1 = 0.73$, $\tau_2 = 0.19$; shares $0.540$ versus $0.342$; $R(1)/R(0) = 0.934$ versus $0.891$; $\hat d_0^{(1)}/\hat d_0^{(2)} = 21.3$. Sweep $\beta = 0.25, 1, 3, 10, 30$ (\code{lr} $= 0.005, 0.02, 0.04, 0.05, 0.05$): mode-1 share $0.314$, $0.328$, $0.409$, $0.540$, $0.581$; mode 2 $0.313$, $0.314$, $0.316$, $0.342$, $0.457$ (at $\beta = 30$ mode 2 condenses too; the window $\hat\delta_1\lesssim1\lesssim\hat\delta_2$ is $\beta\in[2.7, 19]$ here).
  \item S7 ($\ell = 1, 2, 4$): widths read at $\Delta = \ell$: $1.00, 2.00, 4.12$ (mode 1), $1.00, 2.00, 4.07$ (mode 2); $R_1(0) = 4.67, 4.69, 4.76$, $R_2(0) = 0.667$ throughout.
  \item S8 ($k = 0.01, 0.1, 1, 10, 100$ at $\beta = 10$; \code{weight\_lambda} $= 100, 10, 1, 0.1, 0.01$, \code{lr} $= 0.002, 0.02, 0.05, 0.005, 0.0005$ by \eqref{eq:lr-rule}, $\hat c_{\mathrm{data}} = 1.8/\lambda_w$): mode-1 share $0.314$, $0.332$, $0.540$, $0.323$, $0.313$ (peak at $k = 1$); mode 2 $0.313$, $0.314$, $0.342$, $0.322$, $0.313$.
\end{itemize}
The noise-floor smokes of predictions 4 and 5 are deferred with them. Prediction 6 needs $K = 3$; its smoke is the full preset with \code{--set n\_epochs=10000 --set monitor\_int=10}. Read off the member-averaged spectrum, an S1 share has a standard error of $0.083$ at $B = 16$ and the difference of the two $0.117$ (Section~\ref{sec:protocol}(ii)), so at that size the share criteria above are secondary, and the primary smoke criteria are $R(\Delta)/R(0)$ (standard error $0.017$ at $\Delta = 1$); the end-to-end smoke test's $B$ and tolerances are set by W8.

\subsection{Prediction 1: prior-shaped maps (regime a1)}
\label{sec:P1}
\emph{Needs} \code{latent-bounds} (W1), \code{orthonormal-teacher} (W2), \code{whitened-spatial-prior} (W3), \code{linear-theory-module} (W5), \code{spectrum-module} (W6), \code{spectrum-monitor} (W7), all done; \code{linear-ring-presets} (W8) and \code{ring-plots-notebook} (W9); run in W10.
\begin{verbatim}
uv run python -m fct exp init ROOT/p1 --preset linear_ring \
    --axis 'train.beta=[0.25]' --set lr=0.005 --set rate_lambda=0.0
uv run python -m fct exp run ROOT/p1
\end{verbatim}
(one setting, named by the preset's one-value axis; the \code{--axis} repeats the preset's value and may be omitted), and its penalty twin (Section~\ref{sec:presets}). Values: $N = 64$, $P = 256$, $m = n = 8$, $K = 3$, $a = 1$, $\ell = 4$, $k_0 = k_1 = 1$, $\sigma^2 = 0$, $r = 0$, $\beta = 0.25$ ($\delta = 8192$), $2\times10^5$ steps, record every 50, burn-in $10^4$, $B$ members.

\emph{Compute:} per member, the final-state spectra (\code{mode\_spectra}); their average across the $B$ members (row 2, the other contrast, has the same predicted spectrum as row 1 and may be pooled with it); from that average, \code{fold\_q}, \code{spectrum\_summary} of rows 0 and 1 (shares, $R_k(\Delta)/R_k(0)$ for $\Delta\le6$), the ratio row 0 / row 1 and \code{half\_power\_q}, each with a bootstrap over members; \code{predicted\_spectra} for the theory; the burn-in diagnostics of the records; per-member values only for their spread.

\emph{Figure:} three panels: $\hat d^{(1)}_q$, $\hat d^{(2)}_q$ versus $|q|$ (log $y$) with the theory; the ratio versus $|q|$; $R_k(\Delta)/R_k(0)$ versus $\Delta$ with $e^{-\Delta^2/32}$.

\emph{Predict} (limits note, P1, at $\sigma^2 = 0$): $\tau_1 = 0.008$, $\tau_2 < 10^{-3}$; $\hat d^{(k)}_q\approx\lambda_q\hat s_{X,k}^2$; $q_{1/2} = 2$ both (or $3$, $\lambda_3/\lambda_0 = 0.4996$); shares $0.157$ both; $R(1)/R(0) = 0.969$ both; ratio $10.08$, $10.07$, $10.04$, $10.02$, $10.01$, $10.00$ at $|q| = 0..5$ and exactly $10$ beyond ($d^{(2)}_q$ has no floor to reach first); the null rows are zero. Band shares of $|q|\le3$: $0.834$, $0.833$. \emph{Falsified if} $R_k(\Delta)/R_k(0)$ deviates from $e^{-\Delta^2/32}$ by more than $0.05$ for $\Delta\le6$, or the two band shares differ by more than $0.1$, every statistic read off the member-averaged spectrum. (The $0.1$ is this plan's choice, derived from the band shares' standard errors: the difference has standard error $0.041$ at $B = 16$, so $0.1$ is a three-standard-error test from $B\approx24$; Section~\ref{sec:protocol}(ii).) \emph{Secondary} (with $B$ raised as needed): $q_{1/2}$ differing by more than 1, single-$q$ shares by more than $0.03$, the ratio varying by more than $20\%$ over $|q|\le8$ or leaving $10\pm1$ at any $|q|\le18$ (where $\lambda_q/\lambda_0 > 10^{-12}$; beyond, the computed square root of $K$ is at round-off). The $B$ these criteria need is in Section~\ref{sec:protocol}(ii).

\subsection{Prediction 2: rate-penalty whitening (regime b1)}
\label{sec:P2}
\emph{Needs} as P1; run in W10. As P1 with $r = 0.1$: \code{--set rate\_lambda=0.1} ($r = $ \code{rate\_lambda}). Compute and draw as P1, overlaying P1's spectra.

\emph{Predict} (at $\sigma^2 = 0$): $\gamma_1 = 26.7$, $\gamma_2 = 2.7$; $q_{1/2} = 6$ (mode 1) versus $4$; shares $0.077$ versus $0.112$; $R(1)/R(0) = 0.922$ versus $0.953$; the ratio rises monotonically, $1.32$, $1.35$, $1.44$, $1.63$, $2.02$, $2.84$, $4.38$, $6.59$, $8.55$, $9.56$, $9.89$, $9.98$, $10.0$ at $|q| = 0..12$ (no interior maximum: both modes end on the same $\lambda_q$ tail). Band shares of $|q|\le3$: $0.533$ versus $0.710$. \emph{Falsified if} mode 1's $R(1)/R(0)$ or band share is not below mode 2's. \emph{Secondary:} $q^{(1)}_{1/2}\le q^{(2)}_{1/2}$, $\hat d_0^{(1)}/\hat d_0^{(2)} > 3$, or the ratio at $|q| = 9$ below $7$.

\subsection{Prediction 3: target condensation (regime b2)}
\label{sec:P3}
\emph{Needs} as P1; run in W11; the $N = 256$ confirmation needs nothing more.
\begin{verbatim}
uv run python -m fct exp init ROOT/p3 --preset linear_ring \
    --axis 'train.beta=[0.25,1,2,5,10]' --axis 'train.lr=[0.005,0.02,0.04,0.07,0.08]'
\end{verbatim}
This is a $5\times5$ product; run only the diagonal (indices $0, 6, 12, 18, 24$ with \code{--index}; \code{exp list} prints the pairs) or make five ROOTs with \code{--axis 'train.beta=[b]' --set lr=...}; the \code{lr} values are \eqref{eq:lr-rule} with $\hat c_{\mathrm{data}} = 1.04$ ($\beta = 5$ and $10$ are data-limited; a prior-only rule would allow $0.2$ at $\beta = 10$, where the measured $\mathrm{lr}\,\lambda_{\max}$ was $0.16$). $r = 0$. The penalty twin pairs each $\beta$ with its own step $2\beta f\lambda_0/\lambda_w$.

\emph{Compute:} from the member-averaged final-state spectra, the shares and $R(1)/R(0)$ along the sweep, $\hat d_0^{(1)}/\hat d_0^{(2)}$ at $\beta = 10$ and $q_{1/2}$, each with a bootstrap over members. \emph{Figure:} share versus $\beta$ (log $x$) for both modes with the theory; the $\beta = 10$ spectra.

\emph{Predict} (at $\sigma^2 = 0$): at $\beta = 10$ ($\delta = 204.8$): $\tau_1 = 0.83$, $\tau_2 = 0.11$; shares $0.343$ versus $0.165$; $q_{1/2} = 1$ versus $2$; $R(1)/R(0) = 0.985$ versus $0.971$; $\hat d_0^{(1)}/\hat d_0^{(2)} = 53.5$. Sweep: mode 1 $0.157$, $0.165$, $0.188$, $0.281$, $0.343$; mode 2 within $0.157$--$0.165$. Band shares of $|q|\le3$ along the sweep: mode 1 $0.834$, $0.844$, $0.867$, $0.919$, $0.937$, mode 2 $0.833$--$0.844$; $R(1)/R(0)$ of mode 1 $0.969$, $0.971$, $0.974$, $0.982$, $0.985$. \emph{Falsified if} at $\beta = 10$ mode 1's band share is not above mode 2's by at least $0.05$ or its $R(1)/R(0)$ is not above mode 2's, or mode 1's band share at $\beta = 10$ is not above its value at $\beta = 1$ by at least $0.05$. (Both $0.05$ margins are this plan's choice, derived from the band shares' standard errors at $B = 16$: $0.031$ for the mode-1 minus mode-2 difference at $\beta = 10$ (predicted $0.093$), $0.031$ for mode 1 between $\beta = 1$ and $10$ (predicted $0.093$); either difference less $0.05$ is then $1.4$ standard errors at $B = 16$ and three from $B\approx75$.) \emph{Secondary:} single-$q$ shares (mode 1 monotone, mode 2 within $0.03$) and the ratio at $q = 0$ above $20$ at $\beta = 10$.

\emph{Finite $N$:} the share at $\beta = 10$ is $0.172$ at $N = 128$ and $0.086$ at $N = 256$ (Section~\ref{sec:protocol}); a confirmation run at \code{--set N=256 --set spatial.box=256.0} keeps $\ell = 4$ and $\lambda_0$ and must give $R(1)/R(0) = 0.985$, $\hat d_0^{(1)}/\hat d_0^{(2)} = 54$ and the share $0.086\pm0.01$.

\subsection{Prediction 4: white null modes and the floor level (regime c4), deferred}
\label{sec:P4}
\emph{Needs} activity noise at $\sigma^2 = 1$, sampled by one of the two samplers of \path{input_noise_average.tex} (the noise as a $P\times N$ parameter under SGLD, or the closed-form marginal for the linear network); no work item exists, and \code{fct} has no field for $\sigma^2$ (D2). Not run. For when it is: $\sigma^2 = 1$, $r = 1$, $\beta = 0.25$, \code{lr} $= 0.005$. \emph{Compute:} row $K$ (null mean) of the spectrum, its $R(\Delta)/R(0)$; rows 0, 1. \emph{Figure:} the three spectra on one axis with the lines $\sigma_0^2 = 0.5$ and $\sigma^2 = 1$; $R(\Delta)/R(0)$ for $\Delta\le6$. \emph{Predict:} null spectrum flat at $0.5$ with $|R(\Delta)/R(0)| < 0.05$ for $\Delta\ge1$; latent modes $\hat d_0/\sigma_0^2 = 2.0$, $1.9$ and $R(1)/R(0) = 0.18$, $0.14$. \emph{Falsified if} the null $R(1)/R(0)$ exceeds $0.1$ or the null level is off $0.5$ by a factor $1.5$. The level also tells the samplers apart: $0.5$ for the notes' (annealed) measure, $1$ for a fixed (quenched) $\eta$, $0$ for $\eta$ redrawn at every step (\path{input_noise_average.tex}).

\subsection{Prediction 5: noise-dominated (regime c1), deferred}
\label{sec:P5}
\emph{Needs} activity noise, as P4; not run. For when it is: $k_0 = 10^{-3}$ (\code{spatial.input.kernel.amplitude}), $\sigma^2 = 1$, $r = 0$, $\beta = 0.25$, \code{lr} $= 0.005$ ($\rho_1 = 0.27$). \emph{Compute:} $\hat d_0/\hat d_{N/2}$ and $R(1)/R(0)$ per mode. \emph{Predict:} $\hat d_0/\sigma_0^2 = 1.27$, $1.03$; $R(1)/R(0) = 0.025$, $0.003$. \emph{Falsified if} any $R_k(1)/R_k(0) > 0.15$ or $\hat d_0^{(1)}/\hat d_{N/2}^{(1)} > 2$.

\subsection{Prediction 6: no individual-latent cell types}
\label{sec:P6}
\emph{Needs} the runs of P1--P3 with $K = 3$ (W10, W11) and the monitor's \code{hex} and \code{theta\_hist} (\code{spectrum-monitor}, W7, done; \code{spectrum-module}, W6, done); no run of its own; read in W12.

\emph{The statistic.} Each neuron's tuning to the two contrasts, $(t_{i,2}, t_{i,3})$, has the angle $\theta_i = \operatorname{atan2}(t_{i,3}, t_{i,2})$ in the contrast plane $(u_2, u_3)$. At $K = 3$ the three latents project onto that plane as three axes $120^\circ$ apart, so a neuron tuned to one latent, with either sign, sits at one of six angles $60^\circ$ apart. The statistic is the modulus of the hexatic moment, $|m|$ with $m = \langle e^{6i\theta_i}\rangle$ over neurons: $e^{6i\theta}$ takes one value on all six angles, so $|m|\approx1$ when every neuron is an individual-latent cell type, while a rotation-invariant $(t_2, t_3)$, which the theory predicts (independent maps with equal contrast spectra), gives $|m| = O(n^{-1/2})$ over $n$ independent neurons. $|m|$ is frame-free: rotating the frame by $\phi$ multiplies $m$ by a phase $e^{-6i\phi}$. That matters because the contrast singular values are degenerate in expectation ($Pa^2/12$ each), so where the SVD puts $u_2$ within the contrast plane is set by finite-$P$ noise; no choice of frame is needed. The statistic $\langle\cos4\theta\rangle$ of the limits note fails twice: it is measured against that random frame, and it is blind to the effect, since $\cos4\theta$ averages to zero over three axes $120^\circ$ apart (neurons each tuned to exactly one latent give $\langle\cos4\theta\rangle = 0.036$, below its threshold of $0.1$, and $\langle\cos6\theta\rangle = -0.998$). The hexatic moment is meaningful only at $K = 3$: for $K\ge4$, $(u_2, u_3)$ is an arbitrary plane of a larger degenerate contrast space.

\emph{Compute:} per member, $m_b$ of the final state (\code{hexatic\_moment} of rows 1 and 2 of its tuning maps); the complex mean $\bar m$ over the $B$ members (all members share $Z$, hence the frame) and $|\bar m|$, with its error bar from the spread of $m_b$ across members (a bootstrap over members for $|\bar m|$). Neighbouring neurons are correlated ($R(1)/R(0) = 0.97$), so the $64B$ neuron-samples are not independent, and the error bar is taken across members, not as $(64B)^{-1/2}$. Sum \code{theta\_hist} of the members' final states. \emph{Figure:} the summed 36-bin histogram of $\theta$ over neuron-samples, with the uniform density and the six individual-latent angles marked. \emph{Predict:} $|\bar m|\approx0$; \emph{falsified if} $|\bar m| > 0.1$ with $64B\ge10^3$ neuron-samples, i.e.\ $B\ge16$.

\subsection{Prediction 7: correlation length tracks the prior}
\label{sec:P7}
\emph{Needs} as P1; run in W10.
\begin{verbatim}
uv run python -m fct exp init ROOT/p7 --preset linear_ring \
    --axis 'spatial.hidden.kernel.length_scale=[2.0,4.0,8.0]'
\end{verbatim}
(P1's $\beta$ and $r$; the step of the penalty twin depends on $\lambda_0 = \sqrt{2\pi}\,\ell$, so it is one ROOT per $\ell$ there). \emph{Compute:} from the member-averaged spectrum, $R_k(\Delta)/R_k(0)$ for $1\le \Delta\le2\ell$, $R_k(0)$ and \code{gaussian\_width}, each with a bootstrap over members. Per-member widths, if wanted for their spread, need members with $R(\ell)/R(0)$ outside $(0, 1)$ excluded first ($1.8\%$ of members at $\ell = 8$), since \code{gaussian\_width} raises for the whole batch. \emph{Figure:} $R_k(\Delta)/R_k(0)$ versus $\Delta$ for the three $\ell$ with $K(\Delta)$ of \code{wrapped\_gaussian\_profile}; width versus $\ell$. \emph{Predict} (at $\sigma^2 = 0$): widths $2.00, 4.01, 8.06$ (mode 1), $2.00, 4.00, 8.00$ (mode 2); deviations from $K(\Delta)$ below $7.1\times10^{-3}$; $R_1(0) = 26.7, 26.8, 27.2$, $R_2(0) = 2.667$. \emph{Falsified if} a width is off $\ell$ by $20\%$ or $R_k(0)$ changes by $15\%$ across $\ell$.

\subsection{Prediction 8: a common prior scale is non-monotone}
\label{sec:P8}
\emph{Needs} as P1; run in W11. P3's setting ($\beta = 10$, $r = 0$) with $k = 0.01, 0.1, 1, 10, 100$: \code{weight\_lambda} $= 100, 10, 1, 0.1, 0.01$ and, by \eqref{eq:lr-rule} with $\hat c_{\mathrm{data}} = 1.04/\lambda_w$, \code{lr} $= 0.002, 0.02, 0.08, 0.01, 0.001$ (five ROOTs with \code{--axis 'train.beta=[10]' --set weight\_lambda=... --set lr=...}). The data curvature grows with $k$ (measured $\lambda_{\max} = 5.1$ at $k = 10$ and $50$ at $k = 100$, Section~\ref{sec:protocol}(i)), so at $k\ge10$ the step is data-limited and a prior-only rule's $0.2$ diverges at the first step; with the prior initialisation of D1 the prior-only directions need no relaxation, the data-driven ones relax in $10$ to $2\times10^3$ steps, and $2\times10^5$ steps suffice at every $k$ on the whitened path. \emph{Compute:} the two shares of the member-averaged spectrum at each $k$, with a bootstrap over members. \emph{Figure:} share versus $k$ (log $x$), both modes, with the theory. \emph{Predict} (at $\sigma^2 = 0$): mode 1 $0.157$, $0.166$, $0.343$, $0.166$, $0.157$ (peak at $k = 1$); mode 2 $0.157$, $0.157$, $0.165$, $0.165$, $0.157$. Band shares of $|q|\le3$ of mode 1: $0.834$, $0.846$, $0.937$, $0.845$, $0.834$; its $R(1)/R(0)$: $0.969$, $0.971$, $0.985$, $0.971$, $0.969$. \emph{Falsified if} mode 1's band share is monotone in $k$ or its peak is below $0.89$. (The $0.89$ is this plan's choice, derived from the band share's standard error at $k = 1$, $0.014$ at $B = 16$: the predicted peak $0.937$ is $3.4$ standard errors above it; the neighbours at $k = 0.1$, $10$ have $0.028$.) \emph{Secondary:} the same for the single-$q$ share with the peak threshold $0.25$.

\section{Order of work}
\label{sec:order}
Work items in the lab's format (title; type; complexity \code{simple} or \code{complex}; description; inputs; outputs; acceptance criteria), small enough to delegate one at a time; each also updates \path{docs/design.md} and runs the full suite and \code{ruff}. W1--W3 change conventions and come first; W5--W9 build the measurement; W10--W12 are the runs. Each item names its issue slug (\path{docs/issues/<slug>.md}). On \code{sigma0}, W1--W3 and W5--W7 are done and are listed by what they provide; W4 is withdrawn.

\begin{enumerate}[leftmargin=2em,itemsep=5pt]
  \item \textbf{W1. Latents on $[\code{latent\_min}, \code{latent\_max}]$} (\code{latent-bounds}). Type infrastructure; complexity simple. \emph{Done.}\\
  \emph{Provides:} \code{data.latent\_min}, \code{data.latent\_max} (defaults $-1$, $1$, bit for bit the old generator; \code{latent\_min < latent\_max} validated; the replaced field rejected by name), plumbed through \code{LinearTeacherParams}; latents uniform on the interval; \code{latent\_lattice\_dataset} at $\{\min, \mathrm{mid}, \max\}^K$; no data version, since the new field names change every data hash (Section~\ref{sec:D3}).
  \item \textbf{W2. Orthonormal teacher} (\code{orthonormal-teacher}). Type infrastructure; complexity simple. \emph{Done.}\\
  \emph{Provides:} \code{data.teacher: Literal["gaussian", "orthonormal"] = "gaussian"}, dropped from the hashed dump at its default; \code{"orthonormal"} from the \code{Q} factor of the same Gaussian draws, $W^\top W = I_K$; \code{LinearTeacherParams.teacher}.
  \item \textbf{W3. Whitened spatial prior} (\code{whitened-spatial-prior}). Type infrastructure; complexity complex. \emph{Done.}\\
  \emph{Provides:} the field \code{spatial.parametrisation}, \code{"penalty"} or \code{"whitened"} (hashed only as \code{"whitened"}), with $W = S_{\mathrm{rows}}\xi S_{\mathrm{cols}}$ and the plain penalty on $\xi$; \code{net.init = "prior"} (required by \code{"whitened"}; under \code{"penalty"} it draws $W$ from the floored prior), a shared field of \code{ExpConfig}; \code{effective\_weights}; the prior-only step check on both paths and, under \code{"whitened"}, the Lanczos check of $\mathrm{lr}\,\lambda_{\max}(\nabla^2L)$ at run start (\code{largest\_hessian\_eigenvalue}). The parametrisation is an axis of every run of W10 and W11 (Section~\ref{sec:protocol}(v)).
  \item \textbf{W4. Activity noise.} \emph{Withdrawn}: $\sigma\to0$ by policy (Section~\ref{sec:D2}). A sampler of \path{input_noise_average.tex} is the prerequisite of predictions 4 and 5 and is not yet an item.
  \item \textbf{W5. The linear theory from a config} (\code{linear-theory-module}). Type data-science; complexity complex. \emph{Done.}\\
  \emph{Provides:} \path{fct/linear_theory.py} (Section~\ref{sec:D4}): \code{theory\_parameters}, \code{per\_mode\_spectrum}, \code{predicted\_spectra}, \code{predicted\_summary}, \code{ansatz\_free\_kernel}, with $\sigma^2 = 0$ the default and $r = \code{rate\_lambda}$; it reproduces the limits note's P1--P3 numbers from a \code{Config}.
  \item \textbf{W6. Tuning-map spectra: the pure analysis module} (\code{spectrum-module}). Type infrastructure; complexity simple. \emph{Done.}\\
  \emph{Provides:} \path{fct/spectrum.py} (Section~\ref{sec:measure}): the spectra and their summaries, \code{member\_mean\_sem} across members, \code{half\_power\_q} for mean spectra, \code{hexatic\_moment}, and the chain diagnostics \code{record\_sem} and \code{integrated\_autocorr\_time} with their \code{reliable} mask.
  \item \textbf{W7. The spectrum monitor and its wiring} (\code{spectrum-monitor}). Type infrastructure; complexity simple. \emph{Done.}\\
  \emph{Provides:} \code{make\_spectrum\_monitor}, selected by \code{train.monitor\_kind = "spectrum"} (hashed only as \code{"spectrum"}; the shared field \code{monitor\_kind} of an experiment, validated at load); history keys \code{monitor\_spectrum}, \code{monitor\_hex}, \code{monitor\_theta\_hist}, with a member axis for an ensemble; no RNG drawn.
  \item \textbf{W8. Presets and the end-to-end smoke test} (\code{linear-ring-presets}). Type data-science; complexity complex. \emph{In progress.}\\
  \emph{Description:} the two presets and the synthetic-data-first test of the whole pipeline: the smoke ring sampled by SGLD against the theory of W5.\\
  \emph{Inputs:} Sections~\ref{sec:presets}, \ref{sec:smoke}, \ref{sec:protocol}(ii) and (v); W1--W7.\\
  \emph{Outputs:} \path{fct/presets/linear_ring.json} and \path{linear_ring_tiny.json} (grid embeddings, $\sigma^2 = 0$, latents on $[0, 1]$, the orthonormal teacher, \code{init "prior"}, $r = \code{rate\_lambda}$, \code{B} a shared field, $16$ in the tiny preset); a test that runs \code{linear\_ring\_tiny} (S1) through the CLI on a CPU, under both parametrisations, and compares the final-state spectra of the members with \code{predicted\_spectra}.\\
  \emph{Acceptance:} the estimator takes one sample per member, reads every statistic off the member-averaged final-state spectrum with a bootstrap over members, and uses record pooling only as a burn-in diagnostic. The item sets $B$ and the tolerances together for that estimator: on the smoke ring a share has a standard error of $0.083$ at $B = 16$ (Section~\ref{sec:protocol}(ii)). Runtime under two minutes on a CPU; \path{tests/test_presets.py} updated for the two new files.
  \item \textbf{W9. Plots and the results notebook} (\code{ring-plots-notebook}). Type data-science; complexity simple. \emph{To do.}\\
  \emph{Description:} the figures of Section~\ref{sec:measure} and a notebook that shows one setting against its prediction.\\
  \emph{Inputs:} W5, W6, W8; \path{fct/draw.py}; \path{notebooks/build_spatial_exp.py} as the pattern.\\
  \emph{Outputs:} the \code{plot\_*} functions; \path{notebooks/build_linear_ring.py}; a table cell comparing measured and predicted summaries: every statistic of the member-averaged final-state spectrum (shares, $R(\Delta)/R(0)$, widths, ratios, $q_{1/2}$) with its bootstrap error over members, the \code{"penalty"} and \code{"whitened"} settings side by side.\\
  \emph{Acceptance:} the notebook builds and executes on the smoke ROOT of W8; every \code{plot\_*} draws into a given \code{ax} and makes its own figure only without one; tests on synthetic inputs.
  \item \textbf{W10. Run P1, P2, P7.} Type data-science; complexity simple.\\
  \emph{Description:} the weak-target experiments: prior-shaped maps, rate-penalty whitening, correlation length against the kernel width; each under both parametrisations.\\
  \emph{Inputs:} Sections~\ref{sec:P1}, \ref{sec:P2}, \ref{sec:P7}, \ref{sec:protocol}(ii) and (v); W1--W9.\\
  \emph{Outputs:} the ROOTs (\code{p1}, \code{p2}, \code{p7} and their penalty twins), executed notebooks in \code{tmp/}, a short results note stating the falsification verdicts with the measured numbers and their standard errors.\\
  \emph{Acceptance:} every setting \code{done} and its run-start check passed; every compared setting inside the regime its preset documents (Section~\ref{sec:protocol}(vi)), with the preset's exact-versus-predicted numbers quoted; $B$ set per falsifier by Section~\ref{sec:protocol}(ii), at least $16$, and extended with \code{exp extend} for secondary criteria; the burn-in diagnostics reliable on every setting; the verdicts of Sections~\ref{sec:P1}, \ref{sec:P2}, \ref{sec:P7} stated against the predicted numbers, each measured number read off the member-averaged spectrum with its bootstrap standard error over members, and the measured $s_k^2$ of \code{latent\_basis} quoted next to the formula values; the two parametrisations agree on every final-state statistic within their combined standard errors.
  \item \textbf{W11. Run P3 and P8, with the $N = 256$ confirmation.} Type data-science; complexity simple.\\
  \emph{Description:} the strong-target experiments: condensation along the $\beta$ sweep, the non-monotone common scale, and the finite-$N$ check of the $q = 0$ share; each under both parametrisations.\\
  \emph{Inputs:} Sections~\ref{sec:P3}, \ref{sec:P8}, \ref{sec:protocol}(ii), (iv) and (v); W1--W9.\\
  \emph{Outputs:} the ROOTs (\code{p3} with its diagonal, the five \code{p8} roots, the $N = 256$ root, and their penalty twins), executed notebooks, the results note.\\
  \emph{Acceptance:} every setting \code{done}; the regime and $B$ as in W10; the verdicts of Sections~\ref{sec:P3} and \ref{sec:P8} with bootstrap standard errors of the member-averaged statistics; at $N = 256$ the share $0.086\pm0.01$, $R(1)/R(0) = 0.985$ and the $q = 0$ ratio $54$ within $20\%$; the run-start check's $\mathrm{lr}\,\lambda_{\max}$ logged for every $k$ of P8 and below $0.2$; the two parametrisations agree as in W10.
  \item \textbf{W12. Read P6.} Type data-science; complexity simple.\\
  \emph{Description:} the hexatic statistic of prediction 6, read from the final states of the $K = 3$ runs of W10 and W11; no run of its own.\\
  \emph{Inputs:} Section~\ref{sec:P6}; W6, W7; the ROOTs of W10, W11.\\
  \emph{Outputs:} an executed notebook and a section of the results note.\\
  \emph{Acceptance:} $|\bar m|$ with its error bar across members, per setting, and the summed histogram, against the threshold $0.1$ with $B\ge16$.
\end{enumerate}
\appendix
\section{Reproducing the numbers of this note}
\label{app:numbers}
All predicted numbers are from the functions of \path{latex/notes/scripts/check_linear_soln_limits.py} (\code{ring\_spectrum}, \code{latent\_modes}, \code{spectrum\_d}, \code{h\_of}, \code{half\_power\_q}, \code{autocorr\_lag1}, \code{autocorrelation}, \code{wrapped\_gaussian}), called with the parameters stated at each use; the limits note's own numbers are reproduced by the same calls to the digits quoted (checked). The recipe for one parameter set, from the repository root with the script directory on \code{sys.path}:
\begin{verbatim}
import numpy as np, check_linear_soln_limits as L
lam = L.ring_spectrum(N, ell)
# D1 variant: lam = np.maximum(lam, f * lam.max())
sx2, sy2, delta = L.latent_modes(P, K, a, m, k0, k1, n, beta)
# mode k = 0 is the mean direction, k = 1 a contrast
c, d = L.spectrum_d(lam, sx2[k], sy2[k], sig2, r, delta)
tau = sy2[k] * L.h_of(lam, sx2[k], sig2, r).max() / c**2
q_half = L.half_power_q(d); share0 = d[0] / d.sum(); lag1 = L.autocorr_lag1(d)
R = L.autocorrelation(d)
width = np.sqrt(-ell**2 / (2 * np.log(R[int(ell)] / R[0])))
\end{verbatim}
with $(N, \ell, P, K, a, m, n) = (64, 4, 256, 3, 1, 8, 8)$ for the full size and $(16, 2, 32, 2, 1, 4, 4)$ for the smoke, and \code{sig2 = 0} throughout except in the deferred predictions 4 and 5; the ratio is \code{d1 / d2} read at $|q| = 0, 1, \dots$ in FFT order; the $N$-dependence of Section~\ref{sec:protocol} uses $N\in\{64,128,256,512\}$ at $\ell = 4$.

The standard errors and biases of Section~\ref{sec:protocol}(ii) are Monte Carlo draws of the member-averaged estimator under the theory's law: for each repetition, the mean of $B$ members' spectra is \code{d * G} with $G_q\sim\mathrm{Gamma}(B, 1/B)$ for $0 < |q| < N/2$, $G_{-q} = G_q$, and $G_0$, $G_{N/2}\sim\mathrm{Gamma}(B/2, 2/B)$; each statistic (\code{x[0] / x.sum()}, \code{np.fft.ifft(x).real} normalised by its first entry, the width at $\Delta = \ell$, the ratio of the two modes, \code{fct.spectrum.half\_power\_q}) is computed on that mean, and its standard deviation and median minus prediction are taken over $2\times10^4$ repetitions ($10^3$ for $q_{1/2}$), at $B = 16$, $64$, $256$; seed $0$ of \code{numpy.random.default\_rng}. The biases of per-member averaging use one member's draw, $G_q\sim\mathrm{Exp}(1)$ ($\chi^2_1$ at $q = 0, N/2$). The floor's effect in D1 is \code{fct.linear\_theory.solve\_mode} and \code{predicted\_summary} on $\lambda_q^{\mathrm{eff}} = \max(\lambda_q, f\max_q\lambda_q)$ with \code{sig2 = 0}.

The code facts of Sections~\ref{sec:prior} and~\ref{sec:D1} were checked with \code{fct.spatial.build\_spaces} on a spatial config with \code{d=1}, \code{box=64.0}, grid embeddings, hidden \code{length\_scale=4.0}, input and output \code{length\_scale=0.0}, all amplitudes $1$, \code{m = n = 8}, \code{N = 64}, \code{L = 1}:
\begin{verbatim}
spaces = build_spaces(cfg, m=8, N=64, L=1, n_out=8)
# equal exactly; so are box=1.0 with length_scale=4/64:
spaces[0].post.K == L.ring_kernel(64, 4.0)
# equal to np.sort(L.ring_spectrum(64, 4.0)) to 9e-15:
torch.linalg.eigvalsh(spaces[0].post.K)
spaces[0].post is spaces[1].pre                 # True
spaces[0].post.report.n_floored                 # D1's floor counts
torch.linalg.eigvalsh(spaces[0].post.M)[-1]     # D1's lambda_max(M)
\end{verbatim}
for \code{min\_rel\_eig} $\in\{10^{-2}, 10^{-4}, 10^{-6}\}$. The step cost is the wall time of 200 calls of \code{fct.train.compute\_loss} plus \code{backward} on the \code{make\_dataset} output of the full-size \code{DataConfig}, float64, on four threads of the CPU of the machine this note was written on (\code{torch.set\_num\_threads(4)}); the whitened figure is a hand-rolled $S\xi$ forward with the same shapes. Absolute timings depend on the machine; only the order of magnitude is used above.
\end{document}
